\documentclass[11pt]{article}
\usepackage[margin=1.15in]{geometry}
\usepackage{amsmath,amssymb,amsthm}
\usepackage{booktabs}
\usepackage[hidelinks]{hyperref}

\newtheorem{theorem}{Theorem}
\newtheorem{lemma}[theorem]{Lemma}
\newtheorem{proposition}[theorem]{Proposition}
\theoremstyle{remark}
\newtheorem*{remark}{Remark}

\newcommand{\lam}{\lambda}
\newcommand{\eps}{\varepsilon}
\DeclareMathOperator{\supp}{supp}

\title{The one-copy Bernoulli convolution at $2/3$:\\ an $L^2$ dictionary, two theorems, and the routes it closes}
\date{8 September 2026}

\begin{document}
\maketitle

\begin{abstract}
Let $Y=\tfrac13\sum_{k\ge0}B_k(2/3)^k$ with $B_k$ fair independent bits, and let $\mu$ be its law.
Whether $\mu$ is absolutely continuous is a long-standing open question. We record an $L^2$ route: an
elementary reduction, two theorems that make the independence heuristics of the subject exact, an
unconditional partial-sum bound, and proofs that several natural ways of finishing cannot work. Nothing
here settles the question, and we say precisely where each attempt stops. One route is not closed. Recast as a charge on the runs of the exact digit
recursion $2A_{n+1}-3A_n\in\{0,\pm1,\pm2\}$, a strictly local and finite-order combinatorial cost yields
a bound on the octave energy that decays, at $R^{-0.28}$ over nineteen octaves by exact enumeration,
where every previously tested cost grows; making that a theorem requires a large-deviation estimate for
the run statistics of the $\times3/2$ orbit, and we say what is measured and what is not. What survives is a single
statement about $I_K=\int_1^2\prod_{m\le K}\bigl(1+\cos(2\pi x(3/2)^m/3)\bigr)dx$, the integral of a
Riesz product with the non-Pisot ratio $3/2$ over a single octave: membership in $L^2$ is equivalent to
$\sum_K(3/4)^KI_K<\infty$, so what has to be excluded is growth at rate $4/3$ per step rather than mere
unboundedness. Boundedness of $I_K$ is sufficient and strictly stronger. The lattice, the window position and
every Diophantine property of $(3/2)^K$ disappear from the statement; that equivalence is measured here
rather than proved, and we say what its proof would require. Throughout, claims are tested by substituting a biased coin, which leaves the ratio, the lattice
and the period unchanged but makes the measure singular for a small enough bias: an argument that
survives that substitution is proving too much. The measure itself, its characteristic function as an
infinite product, its self-similarity and the functional equation are formalised in Lean~4 with Mathlib,
as are the two theorems and the arithmetic lemma behind them.
\end{abstract}

\section{The reduction to $L^2$}

Write $\mu_n$ for the law of $\{2^nY\}$. For $U\sim\mu_n$ put $2U=H+Z$ with $H=\lfloor 2U\rfloor$,
$u_n(Z)=\mathbb E[(-1)^H\mid Z]$ and $j_n=\log2-H(H\mid Z)$. The $j_n$ are the fibre deficits, and
$\sum_n j_n$ is the relative entropy of $\mu$ against Lebesgue measure, so summability is equivalent to
$\mu$ having a density of finite entropy.

Two remarks. First, $\mu$ is symmetric about $1/2$, so $\mu_n([0,\tfrac12))=\tfrac12$ at every level and
$j_n=I(H_n;Z_n)$ is a mutual information: the question is whether the total information the binary digits
of $Y$ carry about the tail is finite. Second, with $f$ the density and $A=\int f^2$, the inequality
$\log t\le t-1$ gives $\int f\log f-\log A=\int f\log(f/A)\le\int f(f/A-1)=0$, so
\begin{equation}\label{eq:red}
  \sum_n j_n=\int f\log f\;\le\;\log\Bigl(1+2\sum_{k\ge1}|\hat\mu(k)|^2\Bigr).
\end{equation}
\begin{sloppypar}
An $L^2$ density therefore gives summable deficits outright. The target becomes
$S=\sum_{k\ge1}|\hat\mu(k)|^2<\infty$, and numerically $S=0.103155$, so that
$\|f\|_2^2=1+2S=1.206310$.
\end{sloppypar}

\section{Two exact identities}

Since $3Y=B_0+2Y'$ with $Y'$ an independent copy, $\{3Y\}$ and $\{2Y\}$ have the same law, so
\begin{equation}\label{eq:fe}
  \hat\mu(3k)=\hat\mu(2k)\qquad\text{for every }k\in\mathbb Z .
\end{equation}
Writing $k=2^a3^bm$ with $(m,6)=1$ and applying \eqref{eq:fe} $b$ times gives
$\hat\mu(2^a3^bm)=\hat\mu(2^{a+b}m)$: the transform is constant on the anti-diagonals of the $2$--$3$
lattice, whence
\begin{equation}\label{eq:rearr}
  \sum_{k\ge1}|\hat\mu(k)|^2=\sum_{(m,6)=1}\ \sum_{s\ge0}(s+1)\,|\hat\mu(2^sm)|^2 .
\end{equation}
Membership in $L^2$ is thus a lacunary sum with a linear multiplicity weight. Numerically \eqref{eq:fe}
holds to $9\cdot10^{-16}$ for $k\le4000$. One practical consequence: within an octave the dominant
frequencies occur as lattice twins with identical values, for instance $108=2^23^3$ and $72=2^33^2$, so
neighbouring octaves are not independent.

In continuous form $\hat\mu(\xi)=c(\xi)\hat\mu(2\xi/3)$ with $c(\xi)=(1+e^{-2\pi i\xi/3})/2$ and
$|c(\xi)|^2=\cos^2(\pi\xi/3)$. Substituting $\eta=2\xi/3$ in $\int_{\mathbb R}|\hat\mu|^2$ and applying
Plancherel gives, whenever the density lies in $L^2$,
\begin{equation}\label{eq:half}
  \|f\|_2^2=3\int f(x)f(x+\tfrac12)\,dx .
\end{equation}
Measured, $1.206310$ against $3\times0.402103$, agreeing to $6\cdot10^{-7}$; the same relation applied to
the folded measure gives $1.005$ against $1.496$ and fails, so \eqref{eq:half} is not an artefact.
Iterating the functional equation produces a family of such exact relations among finitely many
autocorrelation values, the next being $\tfrac79A(0)=A(\tfrac12)+A(\tfrac34)+\tfrac12A(\tfrac14)$.

Summing \eqref{eq:fe} over $k$ gives one more, on the spectral side. The multiples of $3$ carry
$\sum_{j\ge1}|\hat\mu(3j)|^2=\sum_{j\ge1}|\hat\mu(2j)|^2$, which is the energy on the even indices, so
subtracting from $S$ in two ways,
\[
  \sum_{3\nmid k}|\hat\mu(k)|^2=\sum_{k\ \mathrm{odd}}|\hat\mu(k)|^2 :
\]
the energy on indices coprime to $3$ equals the energy on odd indices. Both sides measure $0.100526$ to
a relative $9\cdot10^{-8}$ at $N=3\cdot10^6$, and the functional equation itself holds termwise to
$1.7\cdot10^{-15}$ over the same range. The rearrangement with the $(s+1)$ weight reproduces
$S=0.103155$ over exactly the $2\cdot10^6$ indices coprime to $3$ below $N$, so the multiplicity is
counted once and not twice, which is the obvious way for this bookkeeping to be wrong. Finally the two
sides of the dictionary meet: $1+2S=1.206310$ against the meet-in-the-middle $\|f\|_2^2=1.206310$,
computed with no shared step.

\section{The seed form and two theorems}

For $\xi\ge3$ there is a unique $K$ with $x=\xi(2/3)^K/3\in[1,3/2)$, and the product splits there:
\begin{equation}\label{eq:seed}
  |\hat\mu(\xi)|=P_K(x)\,T(x)\le P_K(x),\qquad
  P_K(x)=\prod_{j=0}^{K}\bigl|\cos(\pi(3/2)^jx)\bigr|,\quad
  T(x)=\prod_{m\ge1}\bigl|\cos(\pi x(2/3)^m)\bigr| .
\end{equation}
The tail $T$ does not tend to $1$: on $[1,3/2]$ it has supremum $0.046$, mean square $1.11\cdot10^{-3}$,
and a double zero at $x=9/8$. Only the inequality is used below, which is the safe direction; a constant
read off as if \eqref{eq:seed} were an equality is wrong by three orders of magnitude, though the
exponents are unaffected. Equivalently, for integers $k$,
$|\hat\mu(k)|^2=\prod_{i\ge0}\cos^2(\pi k2^i/3^{i+1})$, each factor depending only on $k \bmod 3^{i+1}$.

\begin{lemma}\label{lem:dis}
Let $n_i=2^i3^{K-i}$ for $0\le i\le K$. If $\eps_i\in\{0,\pm1,\pm2\}$ and $\sum_i\eps_in_i=0$, then
$\eps_i=0$ for every $i$.
\end{lemma}

\begin{proof}
Otherwise let $i_1$ be the largest index with $\eps_{i_1}\ne0$. As $\eps_{i_1}$ is prime to $3$, the
$3$-adic valuation of $\eps_{i_1}n_{i_1}$ is exactly $K-i_1$, while every other nonzero term has index
$i<i_1$ and valuation $K-i>K-i_1$. A sum whose terms have a unique minimal valuation has that valuation,
so it is nonzero.
\end{proof}

The frequencies are therefore dissociate; an exhaustive search finds no $\{0,\pm1,\pm2\}$ relation for
$K\le8$.

\begin{theorem}\label{thm:period}
For every $K$,
\begin{equation}\label{eq:period}
  3^{-(K+1)}\sum_{k \bmod 3^{K+1}}\ \prod_{i=0}^{K}\cos^2\!\Bigl(\frac{\pi k2^i}{3^{i+1}}\Bigr)=2^{-(K+1)},
\end{equation}
exactly.
\end{theorem}

\begin{proof}
Expand $\cos^2(\pi\theta)=(1+\cos2\pi\theta)/2$. Averaging over a full period retains exactly those
$\eps\in\{0,\pm1\}^{K+1}$ with $3^{K+1}\mid\sum_i\eps_in_i$; by the valuation count of
Lemma~\ref{lem:dis} that sum has valuation $K-i_1\le K<K+1$ whenever $\eps\ne0$. Only $\eps=0$ survives,
contributing $2^{-(K+1)}$.
\end{proof}

Verified to ten decimals for $K=0,\dots,9$. The independence heuristic behind estimates of this kind is
thus exact at the level of full periods, with no error term.

\begin{remark}
Lemma~\ref{lem:dis}, the orthogonality relation it is combined with, the expansion of the product, and
their assembly are formalised in Lean~4 with Mathlib in \texttt{BC23Full.lean}, which compiles with no
\texttt{sorry}; every theorem there prints only \texttt{propext}, \texttt{Classical.choice} and
\texttt{Quot.sound}. Three further files carry the rest of the chain, also \texttt{sorry}-free with the
same axioms: \texttt{BC23Entropy.lean} for the reduction \eqref{eq:red}, in both its integral and its
finite-sum form; \texttt{BC23PartialSum.lean} for Proposition~\ref{prop:partial}, derived from the period
bound taken as a hypothesis, with the constant $3/2$ shown sharp by a worst-case aggregation check; and
\texttt{BC23Step\-Contraction.lean} for the step and two-level contraction implications.

A further file has the measure itself, which the others do not: \texttt{BC23Measure.lean} builds $\mu$ as
the pushforward of the infinite product of fair Bernoulli measures under $Y$, and derives
$\mu([0,1])=1$, the characteristic function as a limit of finite products, the product formula
$|\hat\mu(\xi)|=\lim_n\bigl|\prod_{k<n}\cos(\pi\xi(2/3)^k/3)\bigr|$, the self-similarity
$\hat\mu(\xi)=\tfrac12\bigl(1+e(-\xi/3)\bigr)\hat\mu(2\xi/3)$, and \eqref{eq:fe} for integer arguments.
Everything before this was stated about an assumed density; these are statements about $\mu$.

The eight files were recompiled as a set on 8 September against the same Mathlib: $39$ certified
statements, no \texttt{sorry}, no \texttt{sorryAx} in any axiom report, no errors, and only
\texttt{propext}, \texttt{Classical.choice} and \texttt{Quot.sound} throughout. As a control that the
check can fail, a copy of \texttt{BC23Measure.lean} with \eqref{eq:fe} altered to
$\hat\mu(3m)=\hat\mu(5m)$ does not typecheck.
\end{remark}

\begin{proposition}\label{prop:partial}
For every $M\ge1$, $\ \sum_{k=1}^{M}|\hat\mu(k)|^2\le\tfrac32\,M^{\log(3/2)/\log3}=\tfrac32M^{0.36907}$.
\end{proposition}

\begin{proof}
Every factor is at most one, so $|\hat\mu(k)|^2\le F_I(k):=\prod_{i=0}^{I}\cos^2(\pi k2^i/3^{i+1})$ for
any $I$. By Theorem~\ref{thm:period} the sum of $F_I$ over a full period $3^{I+1}$ is $(3/2)^{I+1}$.
Take $I+1=\lceil\log_3M\rceil$ and use nonnegativity.
\end{proof}

In correlation-dimension terms this reads $D_2\ge1-\log(3/2)/\log3=0.63093$. It is a power saving over
the $M(\log M)^{-0.46}$ that the proved logarithmic Fourier decay at a rational non-Pisot parameter would
give, and it is comprehensively superseded, which should be said plainly. A two-line folklore argument
already gives $D_2\ge\log2/\log(9/4)=0.85476$; Kleptsyn, Pollicott and Vytnova prove rigorously, by a
computer-assisted method stated for exactly representable rational parameters, that
$D_2(\mu_\lam)\ge0.97580$ for every $\lam\in[0.617,0.743]$, an interval containing $2/3$, with a refined
pass giving $0.9900$ on $[0.5703,0.8)$; and Shmerkin's $L^q$-dimension theorem applies verbatim here,
with exponential separation $\Gamma_k=3^{-(k+1)}$ and an exceptional set containing no rational, giving
$D_2=1$ outright, that is $\sum_{|k|\le R}|\hat\mu(k)|^2=O_\eps(R^\eps)$, though with ineffective
constants. Proposition~\ref{prop:partial} is therefore a short self-contained derivation and nothing
more; it is recorded because it follows from Theorem~\ref{thm:period} in four lines, not because it is
competitive. The three comparisons in this paragraph come from a literature search and I have not read
the primary sources myself; the numbers should be checked against them before being relied on.

\section{What remains, and what is closed}

The octave statement is not new terminology. That $\sum_{R\le k<2R}|\hat\mu(k)|^2\le CR^{-\delta}$ for all
$R$ is, up to $\eps$, Sobolev regularity: $\mu\in L^2_\sigma$ for every $\sigma<\delta/2$, equivalently
Sobolev dimension exceeding $1$ by $\delta$. For a rational non-Pisot parameter nothing of the kind is
published, not even $\delta=0$, and absolute continuity is open for every rational $p/q\in(1/2,1)$ with
$p\ge2$.

It is worth recording exactly why the modern criteria do not reach $2/3$, because the answer is sharper
than ``not yet''. Kittle's theorem (arXiv:2103.12684, Ann.\ Sci.\ ENS 57 (2024) 1191--1230) asks for
$(d\log M_F-h_F)(\log M_F)^2<\tfrac1{27}(\log M_F-\log\lam^{-1})^3\lam^2$. At $2/3$ every input is
exact: the entropy is maximal, $h_F=\log2$, since the digit map is injective; $M_F=M_\lam=3$, the
minimal nonzero $|\nu|$ being $1$; and $\log M_F-h_F=\log(3/2)$ with $\log M_F-\log\lam^{-1}=\log2$. The
left side is $0.4893757$ and the right $0.0054819$, so the criterion fails by a factor $89.27$. That is
not a near miss. The condition is equivalent to $M_\lam<F(\lam)$ with $F$ increasing and bounded by
$2^{27/26}=2.0540$, while every rational $p/q$ in $(1/2,1)$ has $M=q\ge3$: Kittle's criterion can reach
no rational parameter at all. Varju's theorem (arXiv:1602.00261, JAMS 32 (2019)) does admit rationals,
but only for $1-\lam<7\cdot10^{-37}$, against $1-\lam=\tfrac13$ here.

Kittle and Kogler (arXiv:2409.18936, Cor.~1.11) leave the question in an unusually concrete state. Their
condition $(1-\lam)^{-1}h_\lam>C\bigl(\max\{1,\log\log M/h\}\bigr)^2$ collapses at $2/3$, where the
maximum is $1$, to the single requirement $C<3\log2=2.0794$ on one absolute constant; their own
unoptimised estimate is $C\approx10^{10}$. In the other direction the golden ratio, where the measure is
singular, forces $C\ge1.2544$. So absolute continuity at $2/3$ sits inside a factor $1.66$ of one
constant in an existing framework, with about nine of those ten orders of magnitude being machinery
loss rather than mathematics. Provenance, since it matters here: these papers were read and every number
re-derived by two independent passes, but not by me, and I have not read the primary sources myself; the
$10^{10}$ is the authors' own estimate of their own method and is not a theorem, resting as it does on an
unspecified absolute constant from Breuillard--Varj\'u. The claim that no published criterion reaches
$2/3$ is a literature-wide negative built on four papers read and three abstracts, and should be treated
as the current state of a search rather than as a fact.

An octave $[R,2R)$ supplies about $R=(3/2)^K$ consecutive integers, a fraction of order $R^{-1.71}$ of
the period $3^{K+1}$. Through the character expansion,
\begin{equation}\label{eq:err}
  \sum_{k\in W}F_K(k)=2^{-(K+1)}\bigl(L+\mathrm{Err}\bigr),\qquad
  \mathrm{Err}=\sum_{\eps\ne0}c_\eps D(\nu_\eps),
\end{equation}
with $c_\eps=2^{-|\supp\eps|}$, $\nu_\eps=\sum_i\eps_in_i$, $|D(\nu)|\le\min(L,P/2|\nu|)$. The main term
alone gives $E(R)\ll R^{-0.71}$, which suffices, since $S$ is a sum over octaves and
$E(R)\le CR^{-\delta}$ for any $\delta>0$ is enough. Everything turns on $\mathrm{Err}$.

\paragraph{The image, and how the weight sits on it.} The map $\eps\mapsto\nu_\eps$ is injective on
$\{-1,0,1\}^{K+1}$, and its image contains every integer with $|m|\le2^{K+1}-1$, the full range being
$|\nu_\eps|\le\sum_i n_i=3^{K+1}-2^{K+1}$. Uniqueness is the same induction as
Lemma~\ref{lem:dis} read as a digit algorithm: modulo $3$ only the $i=K$ term survives and $\{-1,0,1\}$ is
a complete residue system, so the top digit is forced; divide by $3$ and repeat. Verified exhaustively for
$K\le9$. Writing $c(m)=2^{-\#\{\text{nonzero digits of }m\}}$ for the forced expansion, \eqref{eq:err}
becomes $\mathrm{Err}=2\sum_{m\ge1}c(m)\,\mathrm{Re}\,D_W(m)$, an arithmetically weighted sum with
$\sum_mc(m)=2^{K+1}-1$. That total runs over the whole image, which is spread over
$|m|\le3^{K+1}-2^{K+1}$ and is not consecutive; the consecutive block $|m|\le2^{K+1}-1$ is a sub-block
carrying a shrinking share of the weight, measured by exhaustive enumeration at
$6.11\%,2.74\%,1.23\%$ for $K=8,10,12$ (and $5.20\%,2.31\%,1.03\%$ of the patterns). Injectivity and
the surjectivity onto that block are exact, and were verified by enumeration at the same $K$ with zero
duplicates; what is not true is that the sum runs over consecutive integers, and the shrinking share
above is how far from it the truth is. In the rescaled variable $u=m/2^K$ the window sum has a clean
limit. Since $L/P=1/(3\cdot2^K)$, both
$Lm/P$ and $k_0m/P$ equal $u/3$, so $D_W(m)\sim(3L/2\pi iu)\bigl(e(2u/3)-e(u/3)\bigr)$ and
\[
  \frac{\mathrm{Re}\,D_W(m)}{L}\;\longrightarrow\;g(u):=\frac{3}{\pi u}\cos(\pi u)\sin\!\Bigl(\frac{\pi u}{3}\Bigr),
\]
which matches the exact sum to $2\cdot10^{-9}$ at $m=1$ and to $5\cdot10^{-3}$ at $m=2000$ for $K=10$.
The kernel has mean zero, $\int_0^\infty g=0$ to nine digits, because it is a difference of two sinc
profiles and the window omits the origin; numerically $\sum_mg(m/2^K)=-1.87$ against
$\sum_m|g(m/2^K)|=2147$. What that zero mean gives is an annihilation, not yet an explanation: every
locally polynomial model of the weight $c$ contributes nothing, so the leading term of any smooth model
dies. Whether it also accounts for the \emph{size} of $\mathrm{Err}$ is open and the evidence points both
ways. An independent recomputation puts the smooth-model consequence at $2$--$4\%$ of $\mathrm{Err}$
($-0.0113,-0.0135,-0.0103$ against $-0.3066,-0.3893,-0.4988$ at $K=8,10,12$), which would leave the value
carried by the fine correlation between $c$ and the kernel. Smoothing $c$ over $16$ consecutive
frequencies inside the block instead reproduces the block sum to within $6\%$ at $K=8$ and $0.5\%$ at
$K=12$, and overshoots as the smoothing widens. The two computations run over different index sets, so
they do not yet contradict each other; until that is settled the mechanism should be stated as the
annihilation and nothing more. The sharp form of what remains is
$\sum_{m\ge1}c(m)g(m/2^K)=O(1)$ against $\sum_mc(m)|g(m/2^K)|$ growing geometrically: a local limit
theorem for the weight against one fixed mean-zero kernel. Computing both directly, with $c$ read off by
the digit algorithm, gives for $K=4,\dots,12$ the signed values
$-0.880,-0.930,-0.936,-0.955,-0.955,-0.948,-0.942,-0.930,-0.915$
against absolute values $6.0,9.3,14.0,20.8,30.5,44.3,63.7,91.1,129.5$. The signed sum is the same flat
$-0.93$ found earlier in the sign-pattern coordinates, reached here through a completely different
route, while the absolute sum grows like $1.4^K$. Those signed values are $\mathrm{Err}/L$, not
$\mathrm{Err}$; the deviation is comparable to the main term and cancels most of it, and it is worth
saying so plainly, because read as $\mathrm{Err}=O(1)$ they would settle the problem on the spot.

\begin{sloppypar}
\paragraph{How much cancellation is needed, exactly.} What suffices is $\mathrm{Err}\le CL$: then
$\sum_{k\in W}F_K(k)\le(1+C)2^{-(K+1)}L$, and with $L=R=(3/2)^K$ this is $E(R)\ll R^{-0.7095}$, summable
over octaves with room to spare. So the whole problem is an upper bound on one deviation by a constant
multiple of the main term. The triangle inequality does not give it, and the margin can be measured.
Writing $D(\nu)=\sum_{k\in W}e(k\nu/P)$ exactly and summing $|c_\eps D(\nu_\eps)|$ over all $3^{K+1}$
patterns gives $10.6,\,24.6,\,50.5,\,105.2$ times $L$ at $K=4,6,8,10$, roughly doubling every two
steps, against the true $\mathrm{Err}/L=-0.868,-0.936,-0.955,-0.943$ on the same windows. The identity
itself is exact: $2^{-(K+1)}(L+\mathrm{Err})$ reproduces the direct sum $\sum_{k\in W}F_K(k)$ to every
printed digit at each of those $K$. The phase in $D$ is the whole content --- dropping it and keeping
only $|\sin(\pi L\nu/P)/\sin(\pi\nu/P)|$ turns a quantity near $-0.94L$ into one near $+40L$.
Restricting to the consecutive block flatters the bound but does not rescue it: $6.96,13.4,24.1,43.8$
times $L$ at the same $K$.
\end{sloppypar}

\paragraph{The second moment is exact, and it names the obstruction again.} Average over the window
position instead. Writing $D_{k_0}(\nu)=e(k_0\nu/P)G(\nu)$ with $G(\nu)=\sum_{j<L}e(j\nu/P)$, every
$\nu\ne0$ is killed by the average, so $\mathbb E_{k_0}[\mathrm{Err}]=0$ exactly. For the second moment,
distinct $\eps$ give distinct $\nu_\eps$ with $|\nu_\eps|\le3^{K+1}-2^{K+1}<P$, so no two collide
modulo $P$ and
\[
  \operatorname{Var}_{k_0}(\mathrm{Err})=\sum_{\eps\ne0}c_\eps^2|G(\nu_\eps)|^2 .
\]
The no-collision step is exactly Lemma~\ref{lem:dis}, not the range bound: differences $\nu_\eps-
\nu_{\eps'}$ have coefficients in $\{0,\pm1,\pm2\}$ and can exceed $P$ in absolute value, so the range
alone does not close it. What closes it is that $3^{K+1}$ divides such a sum only if every coefficient
vanishes, which is the same $3$-adic induction, with the same $\{0,\pm1,\pm2\}$ coefficient set the
lemma was stated for.

This reproduces the variance computed directly from every window in a full period, to six digits at
$K=6,8,10,12,14$. Now $\sum_\eps c_\eps^2=\prod_i(1+2\cdot\tfrac14)=(3/2)^{K+1}$ exactly, so the
normalised weight makes the coordinates independent with $\mathbb P(0)=\tfrac23$ and
$\mathbb P(\pm1)=\tfrac16$. Solving $p(1-p)=\tfrac16$ gives $p=\tfrac12(1+3^{-1/2})=0.7886751$, and then
$1-2p(1-p)=\tfrac23$: the auxiliary digit law is exactly $\eta=B-B'$ for $B$ Bernoulli$(p)$, verified to
$3\cdot10^{-17}$ against the product law at $K=4,6,8$.

So the second moment is governed by the correlation integral of the \emph{biased} Bernoulli convolution
at the same ratio $2/3$ with $p=0.78868$. It is tempting to call that a harder instance, on the ground
that the dimension ratio $H(p)/\log(3/2)=1.2719$ is below the unbiased $\log2/\log(3/2)=1.7095$. That
compares the wrong exponent. The quantity governing an $L^2$ statement is
$\log\bigl(p^2+(1-p)^2\bigr)/\log(2/3)$, and at this $p$ the numerator is $\log\tfrac23$ exactly, so the
exponent is $1$ --- the auxiliary measure is exactly $L^2$-critical, which is the same fact as
$\alpha(p)=0$ below, not a harder open problem. For the unbiased measure the two exponents coincide,
both $1.7095$, which is what makes the substitution easy to get wrong. The consequence is favourable
rather than otherwise: a variance linear in $K$ gives a typical $\mathrm{Err}$ of order $L\sqrt K$, and
$R^{-0.70951}\sqrt{\log R}$ is still summable over octaves. Numerically the auxiliary object is well behaved
--- $\operatorname{Var}(\mathrm{Err}/L)$ measures $1.90,2.15,2.54,2.91,3.28$ at $K=6,8,10,12,14$, linear
in $K$ and so logarithmic in $R$, which is exactly why the quantiles of $\mathrm{Err}/L$ over positions
are stable --- but the proof one would need is a harder case of the theorem being proved.

That logarithm is not an accident, and saying why sharpens the whole picture. For the biased measure the
full-period average of the truncated square is $(1-2p(1-p))^{K+1}$, so the main term of the octave sum
has exponent
\[
  \alpha(p)=1+\frac{\log\bigl(1-2p(1-p)\bigr)}{\log(3/2)},
\]
which is negative exactly when $p(1-p)>\tfrac16$, that is for $p$ strictly between the two roots
$\tfrac12(1\mp3^{-1/2})=0.2113249$ and $0.7886751$. But $p(1-p)=\tfrac16$ is precisely the condition
that defined the auxiliary law: the off-diagonal weight of the normalised $c_\eps^2$. So the auxiliary
measure sits exactly on the boundary $\alpha=0$, and a critical exponent produces a logarithm rather
than a power. That is an identity, not a numerical coincidence, and it accounts for the measured linear
growth in $K$. Measured against the prediction, the octave slope is $-0.783$ at $p=\tfrac12$ (against
$-0.7095$), $-0.364$ at $p=0.3$ (against $-0.3435$), $-0.007$ at the critical $p$ (against $0$),
$+0.043$ at $p=0.2$ (against $+0.0488$) and $+0.509$ at $p=0.1$ (against $+0.5106$). At $p=\tfrac12$ the
measured decay beats the main term by about $0.07$, which is the deviation helping rather than hurting,
consistent with $\mathrm{Err}/L\approx-0.94$.

Two thresholds in $p$ therefore sit near each other and should not be confused: the measure is forced
singular below $p=0.14028$, where $H(p)<\log(3/2)$, while the main term stops supplying decay already at
$p=0.2113$. In between, the measure has dimension one and the elementary route gives nothing --- the
same predicament as at $2/3$ itself, but with less room. There is in fact a third threshold between them, and naming it matters for what
follows. The biased factor is $|(1-p)+pe^{i\varphi}|^2=1-2p(1-p)(1-\cos\varphi)$, whose arithmetic average
over an equidistributed phase is $1-2p(1-p)$ --- the quantity $\alpha(p)$ above is built from --- while its
logarithmic average is $2\log(1-p)$, by $\int\log|a+be^{i\varphi}|^2=2\log\max(a,b)$. Annealed and quenched
averages therefore give different thresholds: the typical per-frequency decay turns over at
$(3/2)(1-p)^2=1$, that is $p=0.183503$, strictly between the entropy threshold $0.140277$ and the
main-term threshold $0.211325$. At the latter the annealed log-factor is $-0.40542$ against
$\log(2/3)=-0.40547$, which is $\alpha=0$ to five digits.

This is worth stating because the biased coin is also the one available second control. Every statistic in
this note was calibrated at $1/\varphi$ and $2^{-1/2}$, both algebraic with the singular one Pisot, and
the Pisot reading is forced rather than measured. The biased coin at the \emph{same} ratio $3/2$ supplies
a structurally different singularity --- entropy deficiency rather than arithmetic coincidence --- with
identical frequencies, lattice and window, and its singularity needs nothing deep: the dimension is at
most $H(p)/\log(3/2)<1$ for $p<0.140277$, and a measure on the line of dimension below one is not
absolutely continuous. But it can only be used outside the middle band. Below $0.140277$ the measure is
singular and a criterion should fire; above $0.211325$ the main term supplies decay and it should pass;
between them the measure has dimension one while the main term already diverges, so neither outcome
carries information. A control point placed there has no known answer to calibrate against.

What such a control could establish is detection of entropy-deficient singularity at a non-Pisot ratio,
which is more than this note currently has. What it cannot supply is a measure of dimension one that is
singular for an arithmetic reason --- which is what $\mu$ at the unbiased $2/3$ would have to be, and for
which no example is known. The hole therefore narrows rather than closes.
 That dimension claim deserves its warrant, since the note leans on it more than once. At
$\lam=2/3$ the digit map is injective, and elementarily so: if $\sum_{i\le n}a_i(2/3)^i=0$ with
$a_i\in\{-1,0,1\}$, clearing denominators gives $\sum_ia_i2^i3^{\,n-i}=0$, and reducing modulo $3$ leaves
$a_n2^n\equiv0$, so $a_n=0$ because $a_n\in\{-1,0,1\}$; induction removes the rest. Hence $2/3$ is a root
of no nonzero polynomial with coefficients in $\{-1,0,1\}$, distinct digit strings stay separated at
every level, and exponential separation holds trivially rather than by any Diophantine input. Under that
hypothesis the dimension is $\min\{1,H(p)/\log(3/2)\}$, equal to one exactly when $H(p)\ge\log(3/2)$.
The arithmetic just given is elementary and was checked here; the theorem it feeds is quoted under the
same caveat as the criteria above, from the literature rather than from a reading of its proof. The
point of stating it is the contrast: dimension at $2/3$ is settled, and settled by an injectivity that
costs nothing, while absolute continuity is open --- so the whole difficulty of this note lives in the
gap between full dimension and an $L^2$ density, not in any scarcity of separation.

Applying the substitution to the reduction itself is a check worth making, since a formulation that
quietly hides the conclusion in its error term would pass every other test here. Measuring
$\mathrm{Err}/L$ at the octave position for $K=6,\dots,16$ gives $-0.94$ at $p=\tfrac12$, $-0.78$ at
$p=0.35$, $-0.38$ at the critical $p$, $-0.12$ at $p=0.10$ and $-0.044$ at $p=0.05$, each stable in $K$
to the third digit. It is bounded, and comfortably so, at every $p$ including the singular ones. The
deviation term therefore carries no part of the conclusion: all the decay lives in $\alpha(p)$, and
$\mathrm{Err}=O(L)$ is exactly the supplementary fact it was advertised as. That is the outcome one
wants, and it was not guaranteed.

That the bound cannot be uniform in the position is worth recording separately, since it rules out every
argument that does not look at where the octave sits. Sweeping $k_0$ over a full period, the median of
$\mathrm{Err}/L$ is $-0.34$, the ninetieth percentile $+1.19$ and the ninety-ninth $+6.4$, all stable in
$K$; but the maximum is $13.0,22.7,41.6,75.0,134.4$ at $K=6,\dots,14$, growing like $(4/3)^K$. The
extreme positions are the windows meeting the peak at the origin, where $F_K=1$ and
$\mathrm{Err}\approx2^{K+1}-L$. About $11\%$ of positions have $\mathrm{Err}>L$ at every $K$ tried. The
octave position $k_0=R$ sits near the bottom of the distribution at $-0.94$, and a proof has to know
that.

\paragraph{What the octave position does to the phases.} It does something exact, and it is the
cleanest description of the mechanism we have. The terms that matter are those with $|\nu|\lesssim P/L=
3\cdot2^K$, since $G$ is small outside that range --- which is why a block carrying about one percent of
the weight carries the sum. On that range the phase is $e(k_0\nu/P)$ with $k_0=R=(3/2)^K$, and
\[
  \frac{k_0\nu}{P}=\frac{(3/2)^K\nu}{3^{K+1}}=\frac{\nu}{3\cdot2^K}
  \qquad\text{because}\qquad (3/2)^K\cdot2^K=3^K .
\]
The phase is exactly linear in $\nu$, with slope the reciprocal of the effective range, so it completes
one turn on each side of the origin across precisely the region where $G$ lives. In the rescaled
variable $u=\nu/2^K$ it is the single frequency $e(u/3)$. Since $k_0$ must be an integer the true slope
is $\lfloor(3/2)^K\rceil/3^{K+1}$, but the rounding is immaterial: the ratio to the idealised slope is
$1.00024$ by $K=14$, and substituting the idealised phase reproduces $\mathrm{Err}/L$ to five digits
($-0.93644$ against $-0.93566$ at $K=6$, $-0.91563$ against $-0.91559$ at $K=12$).

Now $g$ is the normalised transform of the octave band, which in this variable is $[\tfrac13,\tfrac23]$,
so $\hat g$ is a multiple of $1_{[1/3,2/3]}(|t|)$. The frequency the octave supplies is exactly
$\tfrac13$: the left endpoint of that band. Evaluating $\hat g$ by tapered quadrature gives
$-0.00000,-0.00000$ at $t=0.20,0.30$ below the band, $1.50000$ at $t=0.40$ and $t=0.50$ inside it,
$-0.00000$ at $t=0.80$ above it, and $0.75000$ at both endpoints $t=\tfrac13$ and $t=\tfrac23$ ---
exactly half the interior value, as a jump discontinuity gives. So the octave, which is precisely the
window whose starting point equals its length, places the phase frequency on the edge of its own band,
where the transform takes half its interior value rather than the full value it would take strictly
inside or the zero it would take strictly outside. That is a clean account of why the answer is a
bounded $O(1)$ number instead of either extreme, and it is specific to $k_0=R$. It is not a proof: the
leading term is only the smooth part of the weight, and the deviation of $c$ from smooth is exactly what
carries the value. But it says what is special about the octave, in one line, exactly.

\paragraph{The octave sits between two exact zeros, and that proves most of the cancellation.} The same
fact has a second face, in the original variable, and this one is an identity with no numerics in it.
Read $F_K(x)=\prod_{i=0}^K\cos^2(\pi x2^i/3^{i+1})$ as a function on the reals. At $x=R=(3/2)^K$ the
factor $i=K-1$ has argument
\[
  \pi\cdot\frac{(3/2)^K\,2^{K-1}}{3^K}=\pi\cdot\frac{2^{K-1}}{2^K}=\frac\pi2 ,
\]
so it vanishes; at $x=2R$ the factor $i=K-2$ has argument $3\pi/2$ and vanishes. The octave $[R,2R]$ is
therefore \emph{exactly} the gap between two consecutive zeros of $F_K$, for every $K$, identically and
not approximately. Confirmed in sixty-digit arithmetic: the identified factor is the smallest of the
$K+1$ at each endpoint and equals zero to working precision at $K=6,10,14$.

The consequence is the strongest structural statement here. Since $k_0\nu/P=\nu/(3\cdot2^K)$ the phase
is a character modulo $3\cdot2^K$, and its $c$-weighted law is exactly uniform: for every $h$ with
$v_2(h)\le K-1$ the expectation $\mathbb E[e(h\theta)]=\prod_n\cos^2(\pi h3^{n-1}/2^n)$ vanishes, since
the single index $n=v_2(h)+1$ puts that argument at exactly $\pi/2$. So replacing $G(\nu)$ by $G(0)=L$
gives $\mathrm{Err}/L=-1$, against the measured $-0.955$: the phases alone annihilate the entire main
term, and the bulk of the cancellation a proof needs is an identity rather than a numerical coincidence.

Three things should be said plainly about how much that buys, and the first of them takes most of it
back. The bound $\mathrm{Err}\ge-L$ is free: $F_K$ is a square, so $\sum_{k\in W}F_K\ge0$, and the window
identity turns that into $L+\mathrm{Err}\ge0$ immediately --- no Fourier expansion, no equidistribution,
at every $p$ and every window position. Measured, $\min_WF_K$ is $8.7\cdot10^{-13}$ at $K=8$ falling to
$7.0\cdot10^{-24}$ at $K=20$, positive throughout. The zeroth-order computation lands exactly on that
free floor: it predicts $A_K=1+\mathrm{Err}/L=0$ where the truth is about $0.05$, so it captures none of
the quantity at issue and bounds it in neither direction. Against $L$ the step looks like most of the
cancellation; against $A_K$, which is what must be bounded, it is none of it.

Second, the zero is at the real point
$(3/2)^K$, while $k_0$ must be an integer; at the integer start the zeroth-order term is small but not
zero, measuring $2^{K+1}F_K(k_0)=4.8\cdot10^{-5},5.5\cdot10^{-5},7.0\cdot10^{-7},2.3\cdot10^{-8}$ at
$K=6,8,10,12$. It is shrinking, so nothing depends on this, but ``exactly'' belongs to the real point.
Third, the character sum obeys
$\sum_\eps c_\eps e(k\nu_\eps/P)=\prod_i\bigl(1+\cos(2\pi kn_i/P)\bigr)=2^{K+1}F_K(k)$, since
$1+\cos2\pi t=2\cos^2\pi t$. So $L+\mathrm{Err}=2^{K+1}\sum_{k\in W}F_K(k)$ identically, and $A_K$ is the
normalised window sum itself --- verified to eight digits by computing it both ways at $K=6,8,10,12$.
So $A_K$ is the same object throughout, and what the exact zero contributes is a clean description of the
octave's position rather than progress against it.

\paragraph{The integrality drops out, and what is left is a Riesz product.} The one reformulation here
that does change the shape of the problem is the passage to the continuum. Comparing $A_K$ with
\[
  \int_1^2\ \prod_{m=0}^{K}\Bigl(1+\cos\bigl(2\pi x(3/2)^m/3\bigr)\Bigr)\,dx
\]
is not an approximation but an identity of integrands. Substituting $k=Rx$,
\[
  \frac{R\,2^i}{3^{i+1}}=\frac{(3/2)^K2^i}{3^{i+1}}=\frac13\Bigl(\frac32\Bigr)^{K-i},
\]
exactly, so $\prod_i\bigl(1+\cos(2\pi k2^i/3^{i+1})\bigr)$ and $\prod_m\bigl(1+\cos(2\pi x(3/2)^m/3)\bigr)$
are the same function with the index reversed. $A_K$ and the integral are therefore the Riemann sum and
the integral of one function, and only quadrature error separates them. Measured, that error is
$A_K-I_K=0.342003\cdot2^{-(K+1)}$, with the constant flat to six digits across $K=9,\dots,27$: it halves
at every step, and the earlier impression that the agreement was erratic --- gaps of $1.7\cdot10^{-5}$ at
$K=8$ against $1.2\cdot10^{-4}$ at $K=12$ --- was noise in an under-resolved quadrature, not a feature.
The true gaps are $6.7\cdot10^{-4}$ and $4.2\cdot10^{-5}$, on the clean geometric ladder. So
$\mathrm{Err}=O(L)$, and therefore the whole $L^2$ target, follows from the boundedness in $K$ of
that integral: a Riesz product with the non-Pisot ratio $3/2$, integrated over a single octave. The
implication is one way only, and the sharper statement is worth having. Writing $I_K$ for that integral
and $B_K$ for the mass of $Q^2dv$ on the block $[(3/2)^K,(3/2)^{K+1})$, the substitution
$\theta=(2/3)^Kv$ relates them, but not by a clean identity: it carries the weight
$\prod_{i\ge1}\cos^2(\pi(2/3)^iv)$, which is bounded above and below on the block and is the reason the
comparison is not exact. Writing $C_K=(3/4)^KI_K/(3B_K)$, the measured values run from $119.0$ at $K=8$ to
$244.3$ at $K=12$, with median $187.1$ over $K=8,\dots,35$, a level shift near $K=11$, and flat at
$0.9968$ per step over $K=13,\dots,35$: bounded within a factor of two, which is all that is needed, and
not a constant. Two earlier drafts got this wrong in opposite directions --- one quoted the ratio at two
endpoints as though it were constant, the second promoted it to an exact identity $B_K=(3/4)^KI_K/2$,
which would force $C_K\equiv2$ and is refuted by the measurements above. With the comparison in hand,
\[
  \mu\in L^2\iff\sum_K(3/4)^KI_K<\infty ,
\]
so $I_K$ may grow, just not at rate $4/3$; boundedness of $I_K$ is sufficient and strictly stronger, since
$B_K=0.9^K$ is summable while $B_K(4/3)^K=1.2^K$ diverges.

What the measurements exclude should be stated without choosing an estimator. Over $K=8,\dots,35$ the
per-step rate of $I_K$ is $1.0032$ as an arithmetic mean of the ratios, $0.9973$ as a geometric mean, and
$0.9829$ by least squares on $\log I_K$; the spread is Jensen's gap, $\sigma^2/2=0.0061$, and the sign of
the drift is not determined by the data. No single rate should be quoted, and earlier drafts quoted two
different ones. What survives every estimator is the exclusion that matters: $\log(4/3)$ sits $13.6$
standard errors above the mean increment and would have produced a factor of $2362$ across the window,
against an observed total factor of $0.930$. So the $L^2$ failure mode is dead over the reachable range.
The resolution is model-dependent and both readings belong here: treating the increments as a random walk
gives a two-standard-error band of $[0.956,1.040]$ on the per-step rate, while under the bounded
hypothesis actually being tested the endpoint estimator resolves a total log-drift of
$0.7506/27=0.0278$. Boundedness itself --- the strictly stronger claim --- is untested either way, and
every sub-exponential law stays invisible: $K^{1/2}$ and $(\log K)^b$ for $b<2.11$ are inside the noise,
while $K^1$ would not be.

The direction the theorem needs is close to rigorous. Since $H$ vanishes at both endpoints of the window,
$|A_K-I_K|\le(2R)^{-1}\int_W|H'|\le\pi D_KI_K$ with $D_K=\sum_i\theta_i\bigl(\int H_i/\int H\bigr)$
measured at $1.10\pm0.03$ for $K=12,\dots,20$. So $\sup_KI_K<\infty$ would give $\sup_KA_K<\infty$ and
hence $\mu\in L^2$, once $D_K=O(1)$ is proved; the converse remains measured only.
No property of $(3/2)^K$ modulo $3^{K+1}$ is available or needed, and the arithmetic of the lattice, the
window position and the choice of integer start have all disappeared. That is the classical form of the
question --- Erd\H{o}s's 1939 setting, and Kahane--Salem territory --- reached from the entropy side.

\paragraph{Why the Riesz-product literature does not reach it, and it is not the ratio.} Two thresholds
are usually run together. Ratio $\ge2$ makes every $\{0,\pm1\}$ combination of the frequencies nonzero,
so the product has mean one; ratio $\ge3$ makes those combinations distinct, so the coefficients are read
off uniquely. Both hold here. The frequencies $\lam_m=(3/2)^m/3=3^{m-1}/2^m$ have strictly decreasing
$2$-adic valuations, so $\sum_m\eps_m\lam_m=0$ with $\eps\in\{0,\pm1\}^{K+1}$ forces $\eps_K$ even, hence
zero, hence all zero by induction --- the same argument as Lemma~\ref{lem:dis} with the roles of $2$ and
$3$ exchanged. The set is dissociate for the coefficients Lemma~
ef{lem:dis} allows, and the constant term of the
product is exactly $1$. It is not dissociate beyond them: $3n_m=2n_{m+1}$ identically for
$n_m=3^m2^{K-m}$, so relations with a coefficient $3$ exist in quantity, and an earlier phrasing here
calling the set dissociate in the strongest sense was wrong.

What fails is not that the frequencies leave $\mathbb Z$ --- scaling by $3\cdot2^K$ returns the integers
$n_m=3^m2^{K-m}$ with ratio $3/2$, a circle Riesz product of the classical kind. It is that the frequency
set depends on $K$: each new factor rescales all the others, so there is no one sequence for the
classical theory to attach to. Since $\lam_0=\tfrac13$ the product has period
$3\cdot2^K$, and $A_K$ is its average over a window that is $1/(3\cdot2^K)$ of that period. The classical
theorems conclude about the mean over the whole group, or about singularity and dimension; none of them
bounds a partial-period average. That last clause is a survey-wide negative, assembled from a literature
search covering Riesz, Zygmund, Peyri\`ere, Brown--Moran, Kilmer--Saeki, Parreau, Ritter and Fan, and I
have not read those sources myself; the dissociateness and constant-term statements around it are
elementary and were checked here in exact arithmetic, but the claim that nobody has bounded such a
quantity should be treated as the current state of a search. On the Bohr compactification, where the classical criterion applies
verbatim, it gives orthogonality to Haar measure, which says nothing about $[1,2]$ because $\mathbb R$ is
Haar-null there. So the gap is not lacunarity but locality.

Two measurements make the position concrete. Over one period at $K=10$ the $3072$ unit windows have mean
exactly $1$, median $0.804$ and maximum $21.4$ at the alignment spike $t=0$ --- so, taking the maximum over unit
windows of the window \emph{mean} of $P_K$ in the $x$ coordinate, $\sup_t\int_t^{t+1}P_K=\|f\|_2^2(4/3)^K$,
measured $1.20646,1.20618,1.20621$ at $K=8,10,12$ against $\|f\|_2^2=1.206310$. A proposed correction
putting that supremum at $2\|f\|_2^2(4/3)^K$ does not survive a direct measurement in this
normalisation; the factor two presumably belongs to a different window length --- the window $[1,2]$ ranks
$486$th, at the sixteenth percentile. That rank is one value of $K$ and it drifts by eight to twenty per
cent over $K=3,\dots,12$, so it locates the octave in the lower part of the distribution without pinning
it. The octave is also the immediate neighbour of the one growing window: the ratio of the window at the
origin to the next is $32,107,269,372,411,451$ at $K=4,\dots,12$, itself of order $(4/3)^K$. And the independence prediction
$\prod_m\int_1^2\bigl(1+\cos2\pi x\lam_m\bigr)dx$ converges to $0.2784$, about five times $A_K$, with the
the shortfall being genuine negative correlation between the factors. Whether it accumulates cannot be
read off this ratio: the denominator converges by $m\approx10$, so the ratio simply tracks $A_K$, and
inherits both the decline through the transient below $K\approx25$ and the flatness after it. Readings
of ``stable'' and of ``monotonically accumulating'' have both been offered here, and each was measured
over a range sitting mostly inside the transient.

The trough has a one-line cause, and it is a single factor. Of the $K+1$ factors, only the lowest is
small on $[1,2]$:
\[
  \int_1^2\Bigl(1+\cos\frac{2\pi x}{3}\Bigr)dx=1-\frac{3\sqrt3}{2\pi}=0.1730067,
\]
against a mean of $1$ over the whole line. It takes the value $\tfrac12$ at both endpoints and vanishes
at the midpoint $x=\tfrac32$. Every other factor averages about one on $[1,2]$ --- measured
$1.000,1.212,1.041,1.154,1.036,1.017$ for $m=1,\dots,6$ --- so the independence value $0.2784$ is
essentially that one integral times a product of terms near unity. The endpoints are themselves zeros of
low-index factors: $x=1$ annihilates the $m=1$ factor and $x=2$ the $m=2$ factor, which is the continuum
face of the exact zeros of $F_K$ at $R$ and $2R$.

So the anatomy of $A_K\approx0.05$ is: a factor $0.173$ from the lowest frequency alone, a factor near
$1.6$ from the rest, and a further factor near $0.18$ from correlation among the factors. The first two
are explicit and bounded. The whole of the difficulty is the third.

That splitting can be made exact, and it turns the target into a correlation inequality with a normaliser
one can write down. The product of the individual means converges absolutely: each factor has
\[
  \Bigl|\int_1^2\bigl(1+\cos2\pi x\lam_m\bigr)dx-1\Bigr|
  =\frac{|\sin4\pi\lam_m-\sin2\pi\lam_m|}{2\pi\lam_m}\le\frac1{\pi\lam_m}=\frac{3}{\pi(3/2)^m},
\]
which is summable, so $\prod_m\int_1^2(1+\cos2\pi x\lam_m)dx$ converges; numerically it is
$0.2783515$, already stable to seven digits by $m=30$. Since that limit is a fixed positive number,
$A_K=O(1)$ is exactly the statement that the ratio
\[
  \frac{\int_1^2\prod_{m\le K}\bigl(1+\cos2\pi x\lam_m\bigr)dx}
       {\prod_{m\le K}\int_1^2\bigl(1+\cos2\pi x\lam_m\bigr)dx}
\]
stays bounded as factors are added --- measured near $0.18$ and flat in $K$. This is a restatement rather
than a bound, but it is the sharpest one available: it names the difficulty as the failure of the factors
to behave independently on a fixed interval, and it supplies the normaliser explicitly, so nothing in the
problem now depends on an unevaluated constant.

Where that dependence comes from is visible at the start of the product. Taking the factors in order, the
ratio above reads $1.0000,\,1.0000,\,0.2652,\,0.3297,\,0.4752,\,0.2683,\,0.2394,\,0.1623$ for the first
one through eight. The first two are \emph{exactly} uncorrelated; the defect is created when the third
joins, and the reason is a coincidence of zeros. Inside $[1,2]$ the lowest factors each vanish once, and
at different places: $m=0$ at the midpoint $x=\tfrac32$, $m=1$ at the left endpoint $x=1$, $m=2$ at the
right endpoint $x=2$, then $m=3$ at $\tfrac43$ and $m=4$ at $1.4815$. Three zeros at the two ends and the
middle leave the product small across the whole interval, which no independence estimate can see, since
each factor separately is small only near its own zero. After that the ratio fluctuates rather than
decaying, later factors contributing correlation and anti-correlation in turn.

Carried to $26$ factors the ratio stays in $[0.1596,0.4752]$, with the maximum attained early, at $n=4$.
Whether it is settling has to be measured carefully: the range over $n\ge n_0$ shrinks as $n_0$ grows for
purely formal reasons, so it is no evidence at all. Fixed-width blocks are not forced to narrow, and over
$n=2$--$6$, $7$--$11$, $12$--$16$, $17$--$21$, $22$--$26$ the block ranges are
$0.236,0.088,0.091,0.039,0.036$ with means $0.316,0.192,0.239,0.180,0.174$. Read alone that looks like
damping toward $0.175$. It should not be read alone: the same quantity carried to $K=40$ has its
five-point running mean fall to $0.041$ and then rise again over the last eight values, so the apparent
settling is inside a window that stops before the next excursion. Taken with the limitation above --- the
oscillation and a logarithm are the same size over every $K$ that can be computed --- the honest summary
is that the ratio has stayed between about $0.16$ and $0.48$ for forty factors and that nothing in the
data decides whether it stays bounded.

One more feature of the parameter deserves recording, because it explains the silence in the literature
better than anything else here. At bias $p$ the same product is $\prod_m(1+a\cos2\pi x\lam_m)$ with
$a=2p(1-p)/(1-2p(1-p))$, and $a=1$ exactly at $p=\tfrac12$. The one classical extension below ratio $3$
requires $\sup_j|a_j|<1$ strictly. It therefore covers every bias except the fair coin, and degenerates
precisely at the case that is open. The endpoint $a=1$ is where this machinery has always stopped.

What is left is exact,
\[
  L\,A_K=\sum_\eps c_\eps\,e(\theta_\eps)\bigl(G(\nu_\eps)-L\bigr),
\]
a decorrelation between the now exactly uniform fractional part of $\theta$ and its own envelope. The
substitution kills this cleanly, as it must: at bias $p$ the vanishing factor becomes $1-4p(1-p)$, zero
only at $p=\tfrac12$. The mechanism belongs to the fair coin, not to the functional equation, and
supplies nothing where the measure is singular. Independently $\mathrm{Err}/L$ has been carried to
$K=38$, a window of $4\,914\,370$ terms, where it reads $-0.954$.

In these terms the remaining task is compact. Writing $A_K=1+\mathrm{Err}/L$, the window identity gives
$E(R)\le2^{-(K+1)}LA_K=\tfrac12R^{-0.7095}A_K$, so $A_K=O(1)$ is the whole of what is left, and the
zeroth-order computation above says $A_K$ has no main term at all. Measured at the octave for every $K$
from $6$ to $40$ --- the last window has $11\,057\,332$ terms --- it runs between $0.0399$ at $K=32$ and
$0.0846$ at $K=12$, a factor of $2.12$, oscillating rather than settling. No trend should be read into
it: log-linear fits give $R^{-0.030}$ from $K\ge6$, $R^{-0.039}$ from $K\ge10$, $R^{-0.011}$ from
$K\ge20$ and $R^{+0.014}$ from $K\ge26$, the sign turning with the fitting window, and the five-point
running mean falls from $0.069$ to $0.041$ and then rises again over the last eight values. An earlier
draft of this note fitted $R^{-0.047}$ on a stride of four and read it as slow decay; that was the
stride's trend, not the sequence's.

Carried further, on an independent code path with composite Gauss--Legendre quadrature and exact integer
phase arithmetic, the integral reaches $K=37$ and settles: over $K=25,\dots,37$ the log-drift is
$+0.00009$ per step at level $0.0427$, and the apparent decay in the earlier ladder is a transient below
$K=25$. Two structural facts come with it, both of which change what a proof should look for. The
recursion is exact, $I_K=I_{K-1}+C_K$ with $C_K=\int_1^2\cos(2\pi f_Kt)P_{K-1}(t)\,dt$ the Fourier
coefficient of the previous product at its own next frequency --- and its multiplier is one, not less:
the geometric mean of $I_K/I_{K-1}$ over $K\ge20$ is $0.9990\pm0.017$. \emph{So the theorem cannot come
from a contraction}, and that route is closed by measurement rather than by argument. What the same data
does show is mean reversion: $\operatorname{Var}(\log I_{K+j}-\log I_K)$ saturates near $0.019$ instead
of growing linearly in $j$, the lag-$8$ variance ratio being $0.28$. Boundedness here is a fluctuation
statement about a driftless positive sequence, not a decay statement.

Which rules out the obvious way to use it. The per-step standard deviation is $0.071$, so independent
increments would give $\operatorname{Var}=0.0403$ over eight steps against the $0.0113$ implied by the
measured ratio; each step contributes about $0.005$ to a sum that therefore diverges linearly. So
$\sum_KC_K^2$ does \emph{not} converge, and a proof cannot come from the increments being square-summable
--- the natural martingale route, and it is closed by the same data that suggests it. Bounded variance
over $j$ steps is bought entirely by cancellation between increments. What has to be proved is that the
$C_K$ are negatively dependent, a mixing statement, and not that they are small.

That target can be localised. Writing $\Phi(u)=\int_1^2\cos(2\pi ut)P_{K-1}(t)\,dt$ and using
$f_{K+1}=\tfrac32f_K$, the product-to-sum identity gives exactly
\[
  C_K=\Phi(f_K),\qquad
  C_{K+1}=\Phi\bigl(\tfrac32f_K\bigr)+\tfrac12\Phi\bigl(\tfrac12f_K\bigr)+\tfrac12\Phi\bigl(\tfrac52f_K\bigr),
\]
all against the same $P_{K-1}$, verified to $10^{-15}$. The last term is negligible, and for a structural
reason rather than by accident: the largest frequency of $P_{K-1}$ is
$\sum_{m\le K-1}(3/2)^m/3=2f_K-\tfrac23$, so $\tfrac12f_K$, $f_K$ and $\tfrac32f_K$ all lie inside the
spectrum while $\tfrac52f_K$ lies outside it. Measured, $|\Phi(\tfrac52f_K)|$ relative to $C_{K+1}$ falls
from $2\cdot10^{-4}$ at $K=8$ to $10^{-10}$ at $K=16$. So the recursion is effectively two-term:
successive increments are one function sampled at $f_K$ and at its two multiplicative neighbours
$\tfrac12f_K$ and $\tfrac32f_K$, and the half-frequency term is the larger of the two in four of the five
$K$ examined. Whether that produces the required anti-correlation is not settled by the five sign pairs
available, three of which flip.

The reduction should not be mistaken for progress. Expanding $\Phi$ over the same characters gives
$\Phi(u)=\tfrac12\sum_\eps c_\eps\bigl[W(\nu_\eps+u)+W(\nu_\eps-u)\bigr]$ with
$W(v)=\int_1^2e(vt)\,dt$ and $|W(v)|\le\min(1,1/\pi|v|)$, so the triangle bound on an increment is again
the weight sitting within $O(1)$ of $\pm u$ --- the same count that defeats the triangle inequality for
$A_K$ itself. Measured, that bound is $21.3,43.4,87.0,172.2$ at $K=8,10,12,14$ against true values
$9.7\cdot10^{-5}$, $5.3\cdot10^{-3}$, $1.5\cdot10^{-2}$, $2.6\cdot10^{-3}$: overshoots of four to five
orders of magnitude, growing. The two-term form is exact and it is a simplification of shape; the
cancellation it needs is the cancellation we started with, in the same amount.

\paragraph{What the closed routes have in common.} Every route closed in this note ends the same way: as
a weighted count of characters near a distinguished frequency, with the true value exponentially below
the count. The triangle bound on $\mathrm{Err}$ over all $3^{K+1}$ patterns, the block weight $(4/3)^K$,
the test-function identity, the second moment, and the increment bound just given are five instances,
with overshoots from a factor $10$ to $10^5$ and growing in $K$.

The common step is not the upper bound; it is \emph{linearisation}. Each of those routes expands
$P_K=\sum_\eps c_\eps e(\nu_\eps t)$ and then bounds term by term, and that discards the cancellation
which is the whole content. But $P_K$ is a product of factors in $[0,2]$ and so is nonnegative: a route
that works with it directly --- covering $[1,2]$ and bounding the function on each piece --- discards
nothing, because for a nonnegative integrand there is no cancellation to lose. That is exactly the
difference between the routes above and the two that remain open. Counting the peaks bounds
$\sum\text{height}\times\text{width}$ for the nonnegative $P_K$; the Erd\H{o}s--Kahane argument is a digit
count over $k\in[R,2R)$ on the object itself. Neither linearises. That is not evidence either will work,
and the peak count walks into $(3/2)^m$ modulo one; but it says which question to ask of a proposed
route, and the question is whether it expands into characters, not whether it ends in an inequality.

\paragraph{The integral is a Birkhoff average for the $\times3/2$ map.} There is a second change of
coordinates, and it is exact. Since $\theta_m(x)=x(3/2)^m/3\bmod1$ satisfies
$\theta_{m+1}=\tfrac32\theta_m\bmod1$, and $1+\cos(2\pi x(3/2)^m/3)=e^{\varphi(\theta_m)}$ with
$\varphi(\theta)=\log(1+\cos2\pi\theta)$, substituting $x=3\theta$ gives
\begin{equation}\label{eq:birkhoff}
  I_K=3\int_{1/3}^{2/3}\exp\Bigl(\sum_{m=0}^{K}\varphi\bigl(T^m\theta\bigr)\Bigr)d\theta,
  \qquad T\theta=\tfrac32\theta\bmod1 .
\end{equation}
Verified against the direct product to all printed digits at $K=6,\dots,16$. So the question is whether a
Laplace average of Birkhoff-type sums along the sequence $(3/2)^m$ modulo one, taken over Lebesgue
measure on a fixed subinterval,
grows at rate exactly one. That is thermodynamic formalism rather than counting: the natural objects are
the transfer operator of $T$ tilted by $\varphi$ and its leading eigenvalue, and the natural statement is
that the growth rate vanishes with a bounded prefactor. It reaches the $\times3/2$ map, as everything
here eventually does, but through pressure rather than through the orbit of a point --- which is the one
door the peak count could not use.

A correction belongs here, because the dynamical reading of \eqref{eq:birkhoff} was wrong and the error
survived several pages of this note. The identity itself is a change of variables and is correct; it was
checked against the direct product to every printed digit. But the sequence $\theta_m$ is \emph{not} the
orbit of $\theta\mapsto\tfrac32\theta\bmod1$. Writing $\theta_m=u_m-k_m$ with $k_m$ the stripped integer,
\[
  \theta_{m+1}=\tfrac32(\theta_m+k_m)\bmod1=\tfrac32\theta_m+\tfrac12[k_m\ \text{odd}]\bmod1,
\]
and the parity of $k_m$ is not a function of $\theta_m$. In exact rational arithmetic $212$ of $429$
consecutive pairs fail the plain map and none fails the skewed form; $x=3$ gives
$\theta=0,\tfrac12,\tfrac14,\tfrac38$, where the plain map sends $0$ to $0$. Nor is this repaired by
another modulus --- reducing modulo $2$ fails as well, and multiplication by $3/2$ is an endomorphism of
$\mathbb R/P\mathbb Z$ only if $\tfrac32k\in\mathbb Z$ for every integer $k$, which holds for no $P$. The
sequence is the orbit of no circle map at all; its natural home is a skew product carrying a parity bit,
or the $2$-adic solenoid. That is precisely the difficulty in Mahler's problem, and it is not something a
one-dimensional transfer operator reaches.

So what \eqref{eq:birkhoff} delivers is a Birkhoff-\emph{type} sum along the sequence
$x(3/2)^m\bmod1$, and not a thermodynamic formalism with an operator behind it. The operator written
below was built for a map that is not the dynamics here, which is the honest explanation of why it
returned a degenerate eigenmeasure concentrated at a fixed point rather than anything about $\mu$: its
leading eigenvalue $4/3$ is a fact about that operator, and the $(4/3)^K$ in this problem is established
independently, by the elementary window decomposition. Everything derived from the closed-form product
rather than from a map is untouched --- the window constant, $\|f\|_2^2=6\int_0^\infty Q^2$, the tail
renewal, the block criterion and the line count all stand, since none of them uses a map.

The way the error survived is worth a sentence. The substitution was verified by checking the
\emph{integrand} in closed form, which is exactly the test that cannot see a false recursion; the
recursion itself was asserted from a one-line manipulation and never checked. It surfaced only because a
separate line of work arrived at a parity-skewed coding and said plainly that it had not proved its
coding equivalent to the one used here.

\paragraph{The atom was a frozen bit, and the annealed rate is exactly one.} The correction above says the
phase is not driven by a circle map; what drives it is $x_{m+1}=\{\tfrac32x_m+\tfrac12b_m\}$ with
$b_m$ the parity of $\lfloor(3/2)^m\theta\rfloor$, a skew product over the two-point fibre. That makes the
earlier atom transparent. For $\theta<(2/3)^{K-1}$ every $b_m$ is zero, which is the one deterministic
branch fixing the origin, and it carries height $2^K$ on measure $(2/3)^K$ --- the whole of the $(4/3)^K$.
Every $\theta$ in the octave leaves that region in a single step and never returns. So the atom is an
artefact of freezing the bit, and the object to look at is the operator averaged over the bits.

There the answer is exact and immediate. For every $x$,
\[
  \bigl(1+\cos2\pi\{\tfrac32x\}\bigr)+\bigl(1+\cos2\pi\{\tfrac32x+\tfrac12\}\bigr)
  =2+\cos3\pi x+\cos(3\pi x+\pi)=2 ,
\]
so averaging over $b_m$ contributes exactly $1$ at every step and the product telescopes:
$\mathbb E_b\bigl[\prod_{m<K}(1+\cos2\pi x_m)\bigr]=1+\cos2\pi x_0$, for every $x_0$ and every $K$. The
annealed octave integral is therefore the $K$-free constant
$\int_{1/3}^{2/3}(1+\cos2\pi\theta)d\theta=\tfrac13-\tfrac{\sqrt3}{2\pi}=0.0576689$, and the annealed
growth rate is exactly one. Nothing is left to prove on that side.

What remains is the whole problem, and one exact rewriting is available. Since
$x_m=\{\tfrac32x_{m-1}+\tfrac12b_{m-1}\}$,
\[
  \cos2\pi x_m=\cos\bigl(3\pi x_{m-1}+\pi b_{m-1}\bigr)=(-1)^{b_{m-1}}\cos3\pi x_{m-1},
\]
so the weight is $w(x_m)=1+(-1)^{b_{m-1}}\cos3\pi x_{m-1}$ identically. Writing $J_K$ for the octave
integral and $A=\tfrac13-\tfrac{\sqrt3}{2\pi}$ for the annealed value, this expresses the increment
$J_m-J_{m-1}$ as an integral of the past product against $(-1)^{b_{m-1}}\cos3\pi x_{m-1}$, and hence
$J_K/A=\prod_{m\le K}(1+\eps_m)$ with $\eps_m$ that increment divided by $J_{m-1}$.

It must be said plainly that the product formula is a telescoping tautology: $\eps_m=J_m/J_{m-1}-1$ by
construction, so verifying $\prod(1+\eps_m)$ to $10^{-16}$ verifies $J_m=J_{m-1}+(J_m-J_{m-1})$, and an
earlier draft here presented it as a result. Its content is the trig identity above, which gives $\eps_m$
an interpretation --- a past-weighted correlation between the parity bit and the phase --- exactly as
$\rho_n$ was given one earlier, and with exactly the same standing: a reformulation, not a reduction. For
the same reason the target $\limsup K^{-1}\sum\log(1+\eps_m)<\log\tfrac43$ is the statement that $J_K$
grows slower than $4/3$ per step, written in other coordinates.

The arithmetic behind the interpretation is real. With $N_m=(3/2)^mv=A_m+x_m$, the relation
$N_{m-1}=\tfrac23N_m$ gives $x_{m-1}=\{((2A_m\bmod3)+2x_m)/3\}$ --- checked here in exact rational
arithmetic, $780$ of $780$ --- so $x_{m-k}$ depends on $x_m$ and $A_m\bmod3^k$ alone while $b_m=A_m\bmod2$.
Since $A_m$ ranges over about $L=(3/2)^m/3$ consecutive integers on the octave, the correlation of $b_m$
with the $k$ preceding phases is $O(3^k/L)$, small for $k$ below
$(m\log\tfrac32-\log6)/\log3=0.36907\,m-1.63$. That is a decay, not an exact vanishing, and the
distinction matters: an independent check finds the inside-window floor falling like $1/L$
($1.1\cdot10^{-4}$ at $m=24$ against $1.5\cdot10^{-5}$ at $m=30$) and the depth growing with $m$ at the
predicted slope, but no step at the boundary --- a factor of $1.9$ between the last inside and first
outside value at $m=24$, and at $m=30$ none at all.

Nor does it license the conclusion an earlier draft drew from it. What is nearly orthogonal to $b_m$ is
the recent-$k$ \emph{truncation} of the product; $\prod_{j<m}w_j$ does not factor into a recent and a deep
piece, so nothing transfers that smallness to $\eps_m$ itself. Measured, the truncated defect inside the
window is about $10^{-4}$ against a full $|\eps_m|$ of $0.048$: the part that is provably small is not the
part that carries the gap.

As for the gap itself, it should not be called a constant. Computed exactly to $K=36$ the ratio runs
between $0.2307$ and $0.4886$, and its median moves from $0.2745$ over $K=8,\dots,35$ to $0.2435$ over
$K=25,\dots,36$. On the octave $|\eps_m|$ is flat across disjoint blocks, $0.0473$ on $m=25,\dots,30$
against $0.0480$ on $m=31,\dots,36$, so it does not decay where it matters; only an ensemble of generic
octaves decays, at $1.4$ standard errors, and this octave is a trough rather than a generic one. What is
established is the absence of drift, $-0.0121$ per step against a critical $\log(4/3)=0.28768$, at about
eighteen standard errors --- with the caveat that twelve increments of a single deterministic curve are
not twelve independent observations, and that the estimate swings by a factor of ten as the window start
moves from $K=8$ to $K=26$. Absence of drift is nevertheless what the theorem needs, since a driftless
walk of step $0.058$ gives $|\log r_K|\sim0.058\sqrt K$.

The margin is wide and is worth stating correctly: the trivial bound $w\le2$ gives $\log2=0.693$ per step
against an allowance of $0.28768$, a factor of $2.41$, while the measurement is $-0.012$. What is asked
for is only an average bound --- no decay rate, no summability, no convergence, no sharp constant. What
is not available is any argument that supplies it, and the deciding measurement is pre-registered: run
this octave alone to $K\approx60$ in exact arithmetic and ask whether $|\eps_m|$ on the disjoint blocks
$m=49,\dots,54$ and $55,\dots,60$ falls below $0.030$. If it does, the ratio is a genuine constant; if it
holds near $0.048$, only the driftlessness is real.

The operator this suggests is explicit, and computing it says what the problem is. Preimages under $T$
are $x=2y/3$ always and $x=2(y+1)/3$ only for $y<\tfrac12$ --- two branches below a half, one above,
which is where the map's difficulty sits --- so with weight $e^\varphi=1+\cos2\pi x$ and $|T'|=3/2$,
\[
  (\mathcal Lg)(y)=\tfrac23\Bigl[\bigl(1+\cos\tfrac{4\pi y}{3}\bigr)g\bigl(\tfrac{2y}3\bigr)
  +\mathbf 1_{y<1/2}\bigl(1+\cos\tfrac{4\pi(y+1)}{3}\bigr)g\bigl(\tfrac{2(y+1)}3\bigr)\Bigr].
\]
Power iteration on grids of $2^{10},\dots,2^{16}$ gives $1.335037,1.335065,1.335071,1.335073$: six digits
of apparent convergence on a number that is not the eigenvalue. The leading eigenvalue is exactly $4/3$,
and the argument is one line. The preimages of $y$ are $2y/3$ and $2(y+1)/3$, so $2y/3=0$ forces $y=0$
while $2(y+1)/3=0$ forces $y=-1\notin[0,1)$; the point $0$ receives from nothing but itself, and the
adjoint's eigen-equation at an atom there reads $\lambda a=\tfrac23(1+\cos0)a=\tfrac43a$. The leading
eigenmeasure is atomic --- the atom at the fixed point together with the fast-decaying chain on the orbit
$(2/3)^n$ into which it leaks --- and a grid cannot carry an atom. The discretisation smears it and
converges, cleanly and consistently, on the wrong value. Four refinements agreeing to six digits was
evidence about resolution, not about the object being resolved. What that atom is has already been
identified above: it is the frozen-bit branch of the skew product, and averaging over the bit removes it
and leaves a rate of exactly one. The operator here is the one with the bit held at zero, so its $4/3$
is correct and describes a single deterministic branch rather than the octave.

The direct measurement settles it and exhibits the mechanism with no spectral theory in it. Split
$\int_0^1\exp\bigl(\sum_{m\le K}\varphi(T^m\theta)\bigr)d\theta$ at $\theta=M(2/3)^K$. On the window,
substitute $\theta=(2/3)^Ku$; since $\theta(3/2)^m=(2/3)^{K-m}u$, the integrand becomes
$2^{K+1}\prod_{j=0}^{K}\cos^2\bigl(\pi(2/3)^ju\bigr)$, so the window equals
\begin{equation}\label{eq:window}
  2\Bigl(\tfrac43\Bigr)^{K}\int_0^M\prod_{j=0}^{K}\cos^2\bigl(\pi(2/3)^ju\bigr)\,du,
\end{equation}
a $K$-free constant times $(4/3)^K$. Measured, the window over $(4/3)^K$ reads $0.402071$, $0.402035$,
$0.402029$, $0.402028$, $0.402028$, $0.402028$, $0.402028$ at $K=12,14,\dots,24$ --- flat to six digits
across twelve steps, where an eigenvalue of $1.335073$ would drift it by $1.6\%$. Everything outside the
window stays bounded over the same range, at $0.32,0.29,0.26,0.23,0.23,0.26,0.27$, and the octave's share
of that remainder is flat between $0.015$ and $0.028$, reproducing from independent code the $A_K$
oscillation measured on the Fourier side.

So the mechanism is elementary. A point with $\theta$ of order $(2/3)^K$ keeps its first $K$ iterates near
zero, where $1+\cos$ is $2$, and so carries height $2^K$; the set of such points has measure $(2/3)^K$;
the product is $(4/3)^K$. That is the whole of the growth, and it is an alignment effect at the fixed
point rather than anything dynamical. The octave is disjoint from the window for every $K$ past $12$, so
its flatness is not cancellation --- it never had this growth to lose. The $4/3$ an earlier draft called
the universal obstruction and then withdrew is real and sits here: branching $2$ against contraction
$2/3$, as measure times height at a fixed point, and not as an obstruction to any particular route.

\paragraph{The growth constant is the norm being bounded.} The constant in \eqref{eq:window} is not a new
number. Write $Q(u)=\prod_{j\ge0}\cos\bigl(\pi(2/3)^ju\bigr)$, so that the window is $2(4/3)^K\int_0^MQ^2$
and $\hat\mu(k)=Q(k/3)$. Since $Q(u)=\hat\nu(u/2)$ for $\nu$ the law of $\sum_{j\ge0}\pm(2/3)^j$, setting
$\rho=\nu*\nu$ gives $Q(u)^2=\hat\rho(u/2)$, whose Fourier transform is $2\rho(2x)$ and is supported in
$|x|\le3$. Poisson summation at spacing $1/3$ then gives, with only $n=0$ landing inside that support,
\begin{equation}\label{eq:constant}
  \|f\|_2^2=\sum_{k\in\mathbb Z}\bigl|\hat\mu(k)\bigr|^2=\sum_{k\in\mathbb Z}Q(k/3)^2
  =3\sum_{n}\widehat{Q^2}(3n)=6\sum_n\rho(6n)=6\rho(0)=6\int_0^\infty Q(u)^2\,du .
\end{equation}
Numerically $6\int_0^MQ^2$ reads $1.205739$, $1.206198$, $1.206295$, $1.206307$, $1.206310$ at
$M=10,10^2,\dots,10^5$, against $1+2S=1.206310$ computed independently on the Fourier side. Hence
$\int_0^1\exp\bigl(\sum_{m\le K}\varphi(T^m\theta)\bigr)\sim\tfrac13\|f\|_2^2\,(4/3)^K$: the constant
standing in front of the operator's leading eigenvalue is the very norm whose finiteness is the question.
The transfer-operator route is circular in the same way the majorant class is --- it returns the quantity
it was meant to bound --- and this is the same fact as the weight profile recorded below, where
$\mathrm{Cum}(u)(4/3)^{-K}\to\frac23\|f\|_2^2u$, reached from the opposite side.

\paragraph{The tail obeys an exact renewal, and $L^2$ is one decorrelation.} $Q$ is self-similar:
$Q(u)=\cos(\pi u)\,Q(2u/3)$, so $Q\bigl(\tfrac32v\bigr)=\cos\bigl(\tfrac{3\pi v}{2}\bigr)Q(v)$ and,
substituting $u=\tfrac32v$ in the tail of \eqref{eq:constant},
\begin{equation}\label{eq:renewal}
  \int_M^\infty Q(u)^2\,du=\tfrac32\int_{2M/3}^\infty\cos^2\bigl(\tfrac{3\pi v}{2}\bigr)Q(v)^2\,dv .
\end{equation}
This is exact. Were $\cos^2(3\pi v/2)$ to decorrelate from $Q(v)^2$ on tail windows it would contribute
its mean $\tfrac12$, the energy would fall by $3/4$ at each $\tfrac32$-step, and the exponent would be
$\log(4/3)/\log(3/2)=0.7095$ --- the endpoint exponent already recorded here. Membership in $L^2$ is
therefore exactly a statement that this one correlation does not persist, which is the phase question in
the only place it can be watched directly across many scales.
Measured on the ladder $N_n=(3/2)^n$ up to $10^5$ in one sweep, the ratio
$\mathrm{corr}_n=\int_{N_n}^{N_{n+1}}\cos^2(3\pi v/2)Q^2\big/\tfrac12\int_{N_n}^{N_{n+1}}Q^2$ has mean
$0.978$ with standard deviation $0.087$ over $n=8,\dots,27$, a standard error of $0.020$, and an
ordinary-least-squares slope of $+0.001$ per step: no drift across twenty scales. The realised block
ratios give exponent $0.775$, faster than the decorrelated $0.7095$ by exactly the amount
$\mathrm{corr}<1$ accounts for, and agreeing with the octave energy $E(R)\approx R^{-0.78}$ measured by
unrelated code.

The power of that null matters more than the null. A correlation biased high by $d$ multiplies the
per-step ratio by $d$ and costs $\log d/\log(3/2)$ of exponent: $d=1.05$ costs $0.12$ and sits $3.6$
standard errors away, while $d=1.35$ --- which would stop the decay outright --- sits at twenty. So the
flat reading is informative rather than merely quiet. What it is not is a proof: $\mathrm{corr}_n$
fluctuates with standard deviation $0.087$ per block, twenty scales of a single curve are not twenty
independent draws, and the permanent limit stands --- over any reachable range, bounded and slowly
divergent cannot be told apart.
Unwound on blocks the renewal gives a threshold, and that is the sharpest form of the question here.
Writing $B_n=\int_{N_n}^{N_{n+1}}Q^2$ and $C_n$ for the same integral against $\cos(3\pi v)$, with
$N_n=(3/2)^n$ and $\rho_n=C_n/B_n$, \eqref{eq:renewal} on a single block reads
\begin{equation}\label{eq:blocks}
  B_n=\tfrac34\bigl(1+\rho_{n-1}\bigr)B_{n-1},\qquad\text{so}\qquad
  B_n=B_0\prod_{k<n}\tfrac34\bigl(1+\rho_k\bigr),
\end{equation}
verified against the measured ladder at all nineteen consecutive pairs to $8\cdot10^{-5}$, which is the
rounding of the tabulated values. Membership in $L^2$ is the summability of that product, and the decay
stops exactly when $\tfrac34(1+\rho)=1$: the critical value is $\rho=1/3$.

That $1/3$ is the same number as the half-shift identity, and it is worth being exact about how little
that says. Summing $V_n=\tfrac43B_{n+1}-B_n$ over all $n\in\mathbb Z$ telescopes, so
$\sum V_n/\sum B_n=1/3$ for \emph{any} positive summable sequence $B_n$: the whole-line ratio is forced
by the renewal alone and carries no information about correlation. It coincides with the critical value
for the same reason --- $1/3$ is where the multiplier $\tfrac34(1+\rho)$ equals one --- so $\rho=1/3$ is
the neutral point of the recursion, pinned at $M=0$ by an identity that reads $A(0)=A(0)$, and free at
every other scale. An earlier draft of this section read the coincidence as the bulk sitting at
criticality; that was a tautology dressed as a fact, and what actually matters is only how $\rho$
behaves in the tail. Measured, it falls fast: the tail cross-ratio reads $0.032,-0.039,-0.008,-0.041$ at
$M=10,10^2,10^3,10^4$, and over the blocks $n=8,\dots,27$ the mean is $-0.0224$ with block standard
deviation $0.0860$ and standard error $0.0192$. The margin to the threshold is $0.356$, about four block
standard deviations. What is not available is any reason it must stay there.
\paragraph{An equivalent reformulation, and what validates it.} Write
$c(n)=B_n^{-1}\int_{N_n}^{N_{n+1}}e^{3\pi iv}Q(v)^2dv$, so that $\rho_n=\operatorname{Re}c(n)$. Then
\eqref{eq:blocks} gives
\[
  \text{if }\ \operatorname{Re}c(n)\le\tfrac13-\delta\ \text{ for all }n\ge n_0,\ \text{ then }\
  B_{n+1}\le\bigl(1-\tfrac34\delta\bigr)B_n,\ \text{ so }\ \sum B_n<\infty,
\]
and $\mu$ has an $L^2$ density. It must be said plainly what this is and is not. By \eqref{eq:blocks},
$\rho_n=\tfrac43B_{n+1}/B_n-1$ identically, so the hypothesis is the conclusion in other coordinates: it
is an equivalent reformulation, not a reduction, and the statement above reads ``if the blocks decay
geometrically then the series converges''. Three further limits belong with it. It is sufficient only
--- $B_n=n^{-2}$ is summable with $\rho_n\to1/3$ and no $\delta$, so a failure of the hypothesis refutes
nothing. Only the real part is needed, so bounding $|c(n)|$ instead is strictly harder than the criterion
it would establish. And it targets $L^2$, which is strictly stronger than absolute continuity, so $\mu$
could be absolutely continuous with $\sum B_n$ divergent.

What survives is worth having anyway, because the reformulation gives $\rho_n$ a meaning the block ratio
does not wear on its face: it is a windowed Fourier coefficient, so ``the blocks decay'' becomes ``the
mass of $Q^2dv$ on $[(3/2)^n,(3/2)^{n+1})$ is not aligned with the phase $v\bmod\tfrac23$''. A windowed
equidistribution statement is a legitimate proof target even when it is equivalent to its conclusion, and
it is a target of a different shape from anything else in this note.

Measured out to $u=1.5\cdot10^6$ --- thirty-four active cosine factors --- $|c(n)|$ has mean $0.112$ over
$n\ge20$ and mean $0.080$ with maximum $0.115$ over $n\ge28$. Reproduced here independently at
$n=28,30,32$: $0.0386$, $0.1151$, $0.0700$, with grid spacings $h=1/16$ and $h=1/48$ agreeing to five
decimals and $B_n$ moving by $10^{-12}$, so the quadrature is converged and the discretisation failure
that produced the spurious $1.335$ above is absent here.

The right way to test a statistic like this is not a standard error but a control at parameters whose
answer is known, and it passes. At $\lam=1/\varphi$, a Pisot number where $\mu$ is singular and
$\sum B_n$ must diverge, the same construction over the same range gives mean $\rho=+0.2404$ against its
own critical value $+0.2361$, with $B_n$ flat to rising and $\rho_n$ stable to eight digits under grid
refinement --- the statistic reads criticality to three digits where criticality is the truth. At
$\lam=2^{-1/2}$, a Garsia number with a bounded density, it gives mean $\rho=-0.4991$ and twenty-six
$e$-folds of decay. At $2/3$ it gives mean $\rho=-0.0231$ with $6.3$ $e$-folds, sitting with the
convergent case and $0.356$ below its own critical value. So the statistic is not blind, and that control
is worth more than any variance estimate computed inside a single deterministic curve.
One limitation outranks the others and should close this section, and it is worse than it first appears.
Every statistic here has been given standing by the same two controls: $\lam=1/\varphi$, where the measure
is singular, and $\lam=2^{-1/2}$, where it has a bounded density. Both are algebraic, and the singular one
is Pisot --- and by Erd\H os's theorem Pisot implies singular, while no non-Pisot singular parameter in
$(1/2,1)$ is known at all. So the controls demonstrate that these statistics detect \emph{Pisot}
criticality and nothing else.

Even that overstates them. At a Pisot $\lam$ the criterion is forced by an identity rather than detecting
anything: the cell at $\xi=1$ has $\|\lam^{-m}\|\to0$ geometrically, so its total cost is bounded while
its width is of order $\lam^{K}$, and that single cell already gives $I_K\ge c\,(2/\lam)^{-K}$, which is
exactly the threshold. The observed near-equality at $\varphi$ --- failing by $0.0037$ --- is therefore a
structural lower bound wearing the costume of a validation, and its flatness is the flatness of one
localised cell rather than of an extensive sum. At $2/3$ the same quantity wobbles by $0.028$ per step,
fifty times larger, because there the sum lives on hundreds of cells. The control tells us nothing about
the measurement at issue.

So if $\mu$ at $2/3$ is singular, it is singular for a reason with no exhibited example, against which the
measured power of every null in this note is not merely unknown but, as the subject stands, unmeasurable.
That $2/3$ sits with the absolutely continuous case is a statement that it is not Pisot-like, which was
never in doubt --- and $2/3$ also sits exactly at the boundary $3\lam=2$, where a smoothed escape rate
equals $\lam$ identically, so at that one parameter a correct computation and a broken one agree.

\paragraph{One register, too small.} The obstructions above were found separately and read as a list; they
are one arithmetic fact. Everything in the product is a function of $\xi$, hence of the pair
$(A_m,x_m)$ with $A_m=\operatorname{round}((3/2)^m\xi)$. The parity bit is $A_m\bmod2$, while the depth-$k$
phase history is determined by $x_m$ together with $A_m\bmod3^k$, since
$x_{m-1}=\{((2A_m\bmod3)+2x_m)/3\}$ and that iterates. So the entire coupling between bit and history is
governed by whether $A_m$ equidistributes modulo $2\cdot3^k$ --- and on the octave $A_m$ ranges over only
about $(3/2)^m$ consecutive integers.

Three consequences follow, and they are the three obstructions. The free window is exactly
$k<m\log(3/2)/\log3=0.36907\,m$, the depth at which $3^k$ still fits inside the available range: that is
the constant measured earlier, now explained rather than observed. The full history wants $3^m$ residues
against $(3/2)^m$ supplied, a shortfall of $2^m$; and $\log3-\log(3/2)=\log2$ exactly, which is
simultaneously the gap between the entropy rate of orbit-sampled digit statistics and that of the leaf
strings the criterion runs over. Against the budget, $\log2/\log(4/3)=2.4094$, so the deficit is not
marginally too large but larger by a factor of two and a half, in every formulation.

One quantity should not be folded into that. Dropping the parity constraint does not merely admit $2^K$
times as many strings: the parity-dropped sum with the exact cost is $(3/2)^K$ identically, a rate of
$0.40547$, so the loss is $0.40547-\Phi\approx0.7255$ rather than $\log2=0.69315$. The extra $0.032$ is the
weight bias of the admitted strings, which carry more than the true ones. The information deficit is
exactly $\log2$; the relaxation penalty is that deficit plus a bias, and calling them one number is tidier
than the arithmetic allows.

What this bounds, and what it does not, is worth being exact about. It bounds any argument that
\emph{decouples} the bit from the deep history: such an argument asks $A_m$ to equidistribute past its own
range and pays $\log2$ per step, whatever its cost model, cost order or automaton size. It says nothing
against an argument that keeps them coupled and estimates the actual correlation, whose measured
per-step defect is about $0.048$ --- far below what any worst-case relaxation must pay. That remaining
route is the decoupling statement for a leaf of the $2$-adic solenoid, and the deficit explains precisely
why it cannot be reached by relaxation, only by genuine equidistribution input about $(3/2)^m$.

Two structural cautions belong beside it, both established here on a second operator. The killed transfer
operator for $x\mapsto\lam x+\eta$ on its trapping window has escape rate exactly $\lam$ once
$3\lam\le2$, because each branch preimage has length $4\lam$ and fits inside the window; so any smoothed
version of the problem returns the critical coincidence by arithmetic, at every resolution rather than in
a limit. An exact-overlap discretisation confirms it, giving $\lam$ flat to six digits at grids of $500$
through $32000$ --- and a floating-point particle filter converges confidently to
$0.618014\pm0.000039$, which is $\lam$ to five places, erasing the $0.37\%$ signal it was built to see,
while the exact filter recovers $0.6202986$. Smoothing does not degrade the answer gradually; it replaces
it. That is the same failure as the spurious $1.335073$ above, in a different operator, now with an
analytic reason behind it. (An earlier draft reported a grid ladder drifting away from the truth; that
does not reproduce, and the correct statement is the stronger one just given.)

There is a second reason to read the escape rate at $2/3$ narrowly. That parameter sits exactly at the
boundary $3\lam=2$, where the smoothed rate equals $\lam$ identically, so a correct computation and a
broken one return the same value there; and $\rho=\lam$ is the first-moment count already known to be
exactly critical, since the cells per unit and the mean cell length multiply to one at every level. A
Pisot-sized excess is therefore excluded at $2/3$ by a statistic validated on Pisot parameters, at a
place where Pisot behaviour was already ruled out by the dissociativity above. It is a consistent
reading, not an informative one.

Its limit is exact and should be stated as such. The Pisot control has no transient --- $\rho_n$ is
essentially constant from the first block --- so it establishes that the statistic detects criticality
that is present from the start, and says nothing about a measure which mimics decorrelation for thirty
scales and then re-correlates. What the measurement excludes is therefore the onset of re-correlation
before $n\approx34$; onset after that is untouched, and no finite ladder of blocks can touch it. An
earlier draft of this section reported a drift in the segment means and exponents as the signature such a
tail would leave. That is withdrawn: the segments were nested, one block carries about sixty per cent of
the apparent gap, and on disjoint windows there is no gradient at all.
There is a structural reason $|c(n)|$ is small, and it is the same arithmetic that runs through the rest
of this note. Rescale the block by $v=(3/2)^nw$ with $w\in[1,3/2)$. Since
$Q\bigl((3/2)^nw\bigr)=\prod_{i=1}^{n}\cos\bigl(\pi(3/2)^iw\bigr)\,Q(w)$, the integrand becomes a Riesz
product of ratio $3/2$,
\[
  \prod_{i=1}^{n}\cos^2\bigl(\pi(3/2)^iw\bigr)
  =\prod_{i=1}^{n}\tfrac12\bigl(1+\cos(2\pi(3/2)^iw)\bigr),
\]
weighted by $Q(w)^2$, and $c(n)$ is its normalised Fourier coefficient at frequency $(3/2)^{n+1}$. That
frequency is not in the spectrum. The coefficients live on $\sum_{i\le n}\eps_i(3/2)^i$ with
$\eps_i\in\{0,\pm1\}$; were $(3/2)^{n+1}$ such a sum, clearing denominators would give
$3^{n+1}=\sum_{i\le n}\eps_i3^i2^{\,n+1-i}$, and reducing modulo $3$ kills every term but $i=0$, forcing
$\eps_0=0$; dividing by $3$ reproduces the same equation one degree lower, so every $\eps_i$ vanishes and
$3^{n+1}=0$. An exhaustive search over $\{0,\pm1\}^{n+1}$ confirms it for $n\le9$.

So $c(n)$ would be exactly zero on a full period, and what the block measures is only the leakage from
the finite window $[1,3/2)$ and the weight $Q(w)^2$. That is why the measured values sit near $0.08$
rather than near $1$. It is an explanation and not a bound, and the gap between the two is the one this
note keeps arriving at: expanding the product into its $3^{n+1}$ characters and estimating the window
transform against each overshoots by exactly $2^n$, since the coefficients have total variation
$\prod(1+\tfrac12+\tfrac12)=2^{n+1}$ against a mean of $1$. Nor does exactness help by itself --- the
distance from $(3/2)^{n+1}$ to the nearest attained sum is
$|3^{n+1}-\sum\eps_i3^i2^{\,n+1-i}|/2^{\,n+1}$, whose numerator is a nonzero integer and so at least one,
which bounds the distance below only by $2^{-(n+1)}$. Lines may therefore approach the target
exponentially closely, and a proof must count how many do, not merely observe that none is exact. That is
the same distinction already drawn for the dissociateness lemma above.
That count can be done exactly, and it settles the matter against this line of attack. Writing
$S_\eps=\sum_{i\le n}\eps_i3^i2^{\,n+1-i}$, the number of $\eps\in\{0,\pm1\}^{n+1}$ with
$|3^{n+1}-S_\eps|\le D\,2^{\,n+1}$ --- that is, with frequency distance at most $D$ --- is, by
meet-in-the-middle at every $n$ from $4$ to $16$ and every $D\in\{\tfrac12,1,2,4\}$, exactly
\[
  N(n,D)=D\,2^{\,n},
\]
with the smallest attained integer gap equal to $1$ throughout, so the nearest line lies at frequency
distance $2^{-(n+1)}$. And $D2^n$ is precisely what uniform spreading predicts: $3^{n+1}$ values spread
over a range of about $4\cdot3^{n+1}$ have density $\tfrac14$ per integer, giving
$2D2^{\,n+1}/4=D2^{\,n}$. The measured count matches to the last digit at every $n$.

So there is no arithmetic conspiracy near this frequency and no Diophantine obstruction to be found; the
sum-set is as evenly spread as a random set of the same size, and the difficulty is density rather than
arithmetic. The consequence is negative and worth stating plainly: counting cannot prove the criterion.
Any estimate that treats the leakage line by line in absolute value sums $2^n$ terms of size $O(1)$ and
fails by exactly the factor that closes every linearising route in this note --- the same $2^n$, reached
now from a third direction. What the criterion needs is cancellation among the leakage phases. That
leaves it a genuine reformulation of the question rather than a way around it, which is the honest
description of everything in this section.
\paragraph{Erd\H os--Kahane at a fixed parameter: the count is an identity, but the route is live.} The one remaining route
that does not linearise survives, but not in the form it is usually given: its counting step is empty here,
and what carries any decay is the per-cell charge instead. The apparatus is exact arithmetic here. With $x_n=(3/2)^{n-N}\xi$,
$A_n=\operatorname{round}(x_n)$ and $d_n=x_n-A_n$,
\begin{equation}\label{eq:ek}
  e_n:=2A_{n+1}-3A_n=3d_n-2d_{n+1}\in\{0,\pm1,\pm2\},\qquad e_n\equiv A_n\bmod 2,
\end{equation}
checked in rational arithmetic over $12{,}000$ steps without a failure, so the digit string of $\xi$ is
the itinerary of a full-branch expanding map of slope $3/2$ and $\xi$ is its initial condition.

The engine of the method is that the number of admissible strings times the parameter measure carried by
each tends to zero. Over $\lam$ that measure is exponentially small. At a fixed $\lam$ it is not: the
number of cells per unit of $\xi$ is $\sum_m(2/3)^m=3$ exactly, independent of $N$, the mean cell length
is $1/3$, and cylinder measure times contraction is $(3/2)^N(2/3)^N=1$. Three independent derivations
agree, so the product is exactly $1$ and the counting step does no work at any $N$. The mechanism is that
$\xi$ enters its own orbit through a single linear coordinate, $dx_n/d\xi=(3/2)^{n-N}\le1$, so the
level-$n$ constraints are nested and intersect to the tightest one; over $\lam$ they are transversal and
multiply. Rationality of $1/\lam=3/2$ is what supplies the exact finite alphabet in \eqref{eq:ek} and
simultaneously empties the count.

Two numerical facts are worth keeping. Membership in $L^2$ needs a cost exceeding
$\tfrac12\log(3/2)=0.20273$ nats per step; the support-only version of the method delivers $0.1832$, and
a version aware of $|e_n|$ delivers a mean of $0.304$ --- half again what is required. And yet the bound
it actually yields, $R\,\mathbb E[e^{-2\,\mathrm{cost}}]$, \emph{grows}, as $R^{+0.126\pm0.027}$; the
support-only version grows as $R^{+0.34}$; the truth decays as $R^{-0.8}$. Using $e^{-2\mathbb E[\cdot]}$
in place of $\mathbb E[e^{-2\cdot}]$ is Jensen's inequality in the wrong direction, and here it reverses
the sign of the conclusion rather than merely weakening it: the annealed average is dominated by the rare
strings that avoid $|e_n|=2$ entirely. It is worth recording as a general caution, since the mean cost
comfortably clears the target while the method fails.

What that does \emph{not} establish is that the route is closed, and an earlier draft of this paragraph
said it was. The claim was that any cost function good enough to decay must encode the cylinder geometry
and would therefore be evaluating the integral rather than counting strings. Exact cell enumeration
refutes it. Charging each maximal run of nonzero $e_n$ its own constrained optimum --- a relaxation, since
dropping the couplings between runs only raises the maximum and so lowers the charge, giving a valid
majorant, audited directly at $+1.1\cdot10^{-13}$ over all $82{,}738$ runs at $N=21$ --- produces a bound
that \emph{decays}, at $R^{-0.28}$ to $R^{-0.31}$, falling monotonically at every $N$ from $11$ to $30$
and by a factor of $9.1$ over nineteen octaves. On the identical machinery the per-symbol cost grows at
$R^{+0.11}$, so the test could plainly have come out the other way; and the exact enumeration removes the
sampling caveats that attended the earlier Monte Carlo estimate of the same growth.

The crossover sits at order two. Charging only disjoint blocks of at most two consecutive jump symbols
already decays, at $R^{-0.15}$, and order six is within two per cent of the full cost. That cost is a
strictly local, finite-order, purely combinatorial function of the digit string: it discards every
coupling between runs and everything at the roughly one third of indices carrying $e_n=0$, and at $N=25$
it sits some $3\cdot10^5$ above the true octave energy while decaying at less than half the true rate.
It is not the integral in disguise.

So the route is live, and the obstruction is not where it appeared to be. It is not the counting step ---
that really is an identity, and no variant will draw decay from it --- but the free energy of the digit
process. What a theorem needs here is a rigorous large-deviation estimate for the run-pattern statistics
of the actual $\times3/2$ orbit, and that process is demonstrably not independent: an independent model
misses the measured per-symbol exponent by about $0.05$ nats per step, in the direction and of the size
that the measured over-dispersion of the $|e_n|=2$ counts predicts. A numerically decaying bound over
nineteen octaves is a measurement, not a proof, and these are nested octaves rather than independent
observations, so what carries the sign is the strict monotonicity rather than any standard error.



\paragraph{The recursion closes on the integers, and that is why it cannot help.} Write
$T(f)=\int_{\mathbb R}\cos(2\pi fv)Q(v)^2dv$. The self-similarity $Q(v)^2=\cos^2(\pi v)Q(2v/3)^2$ together
with $\cos^2(\pi v)=\tfrac12+\tfrac12\cos2\pi v$ gives an exact three-term relation,
\begin{equation}\label{eq:trec}
  T(f)=\tfrac34T\bigl(\tfrac{3f}{2}\bigr)+\tfrac38T\bigl(\tfrac{3(f+1)}{2}\bigr)
       +\tfrac38T\bigl(\tfrac{3(f-1)}{2}\bigr).
\end{equation}
Since $T(f)=2\rho(2f)$ and $\rho=\nu*\nu$ is supported in $|x|\le6$, $T$ vanishes for $|f|>3$; the
three-term step at an integer reaches only $0,\tfrac32,3,\tfrac92$, and that truncation closes the
system. Taking $f=0,1,2,3$ in turn gives $T(3/2)=T(0)/3$, then $T(1)=\tfrac58T(0)$, then
$T(2)=\tfrac18T(0)$, with $T(3)=0$ consistent. Measured on the direct integral to $M=10^5$:
$0.62500006$, $0.33333349$, $0.12500006$, and $T(3),T(7/2),T(4)$ zero to $10^{-15}$. The points that do
not close look nothing like rationals, $T(1/2)/T(0)=0.88359189$ and $T(5/2)/T(0)=0.02328172$, which is
what the escape into quarters predicts.

In terms of the autocorrelation itself, with $A(s)=3T(3s)$, these read
$A(1/3)=\tfrac58A(0)$, $A(1/2)=\tfrac13A(0)$, $A(2/3)=\tfrac18A(0)$ and $A(1)=0$. The middle one is the
half-shift identity; the outer two came from the same recursion. The last is free, since $f$ lives on
$[0,1]$, and so serves as a check on the correspondence rather than as a result. A second check needs no
computation at all: $A$ is even, supported in $[-1,1]$, and integrates to $(\int f)^2=1$, so
$\int_0^1A=\tfrac12$ exactly, and the four predicted values $1.206310,0.753944,0.150789,0$ give
$0.502629$ under the trapezoid rule at $h=1/3$ --- inside its own error on four points, and overshooting
as it must for a convex $A$. All four are confirmed at once by a constraint none of them was derived
from.

None of it can settle the question, and the reason is structural rather than technical.
Equation~\eqref{eq:trec} is linear and homogeneous in $T$, so it determines $T$ up to a single scalar,
and that scalar is $T(0)=\|f\|_2^2/3$ --- the unknown. Every identity in this family constrains the shape
of the autocorrelation and none constrains its scale; and each presupposes what it would prove, since
were $\|f\|_2^2$ infinite the relations would read $\infty=\infty$. The measure's defining property is
self-similarity, self-similarity generates homogeneous linear relations, and finiteness of the $L^2$ norm
is exactly the inhomogeneous information that self-similarity cannot supply. That is the same fact as
the observation running through this note, that every closed route ends in a phase count: the phases are
what the functional equation does not carry. A proof must inject something that is not a linear
identity --- the arithmetic of $(3/2)^m$, or a counting argument on the nonnegative product.




A numerical caution belongs with \eqref{eq:birkhoff}, since it cost an hour. Generating the orbit by
iterating $\theta\mapsto\tfrac32\theta\bmod1$ in floating point multiplies the error by $3/2$ at every
step and destroys the fractional part: done that way the right-hand side reads $0.045,0.097,0.186,0.341,
0.614,1.100$ at $K=6,\dots,16$, growing, against a true value near $0.05$. Computing each
$\theta_m=\theta(3/2)^m\bmod1$ in closed form instead reproduces the direct value exactly. A chaotic map
must not be iterated numerically when its orbit is available in closed form.

\paragraph{What the peak count would have to supply.} Being outside the shape is not the same as being
easy, and the requirement can be written down. Cover $[1,2]$, bound $P_K$ on each piece, and sum height
against measure. A point meeting a fraction $\alpha$ of the $K$ conditions $\|t\lam_m\|<\delta$ carries
height about $2^{\alpha K}$; if those conditions behaved independently, with $p_0=2\delta$, the measure
of the set meeting at least $\alpha K$ of them would be $e^{-KD(\alpha\|p_0)}$, and the level would
contribute $\bigl(2^\alpha e^{-D(\alpha\|p_0)}\bigr)^K$. Optimising over $\alpha$ gives the log moment
generating function exactly,
\[
  \sup_\alpha\bigl[\alpha\log2-D(\alpha\|p_0)\bigr]=\log(1+p_0),
\]
so the independence estimate is $(1+2\delta)^K$ --- which is just the first moment
$\mathbb E\,2^{\#\text{met}}$, and exceeds one for every threshold. Numerically $1.20^K$, $1.30^K$,
$1.40^K$ at $\delta=0.10,0.15,0.20$. No choice of $\delta$ rescues it.

Against that, the measured carrier has measure $0.7119^K$ and height $1.4046^K$, whose product is $1$ by
construction, being the flatness of $I_K$ rather than an independent fact. The gap between $0.7119$ and
the independence values $0.80$, $0.9165$, $0.9798$ at $\alpha=\tfrac12$ is what a proof would have to
supply: the conditions $\|t\lam_m\|<\delta$ must be shown \emph{positively correlated at a definite
rate}, not merely counted. That is the honest price of the route that survives, and it is why it runs
into the distribution of $(3/2)^m$ rather than around it.

A caution applies to all of these increment statistics, and to several figures earlier in this note.
There is a transient below $K\approx25$, and the range a laptop reaches in reasonable time lies wholly
inside it. Measured here, the log-drift of $I_K$ per step is $-0.0165$ over $K=7,\dots,26$ and $-0.0247$
over $K=15,\dots,26$, against $+0.00009$ over $K=25,\dots,37$ from the longer run: two regimes, not a
disagreement. The same split explains the standard deviation of the relative increments, $0.089$ to
$0.137$ below $K=26$ against $0.071$ above $K=20$. It is also the common cause of the trend readings
withdrawn above --- the stride fit, the sign-turning window fits, the declining correlation ratio, and
the dispute over whether the correlation accumulates. None of those was a property of the object. Any
statement here fitted below $K=25$ describes the transient.

The mass has a fixed carrier. At $K=20,24,28$ the same peak locations recur --- $1.869287$, $1.878192$,
$1.871758$, $1.881893$, $1.622273$ --- with heights growing like $1.40^K$ and widths like $(2/3)^K$,
and the set is closed under $t\mapsto2t/3$, the ratio-$3/2$ dynamics acting on it. The top $10^{-3}$ of
Lebesgue measure carries $55\%$ of the integral at $K=20$, $80\%$ at $K=26$ and $95\%$ at $K=32$. So
$A_K=O(1)$ would follow from counting those peaks and bounding height against width, which is a concrete
programme rather than a hope. These measurements are an agent's, on its own code, and their adversarial
check had not returned when this was written.

They do cohere, and the check is not circular, because one of the four numbers is forced rather than
measured. Height $1.4046^K$ against width $(2/3)^K$ gives $0.9364^K$ per peak, so a flat integral
requires the number of peaks to grow like $1.0679^K$; that count was never measured. It then predicts a
carrying set of measure $(2/3\cdot1.0679)^K=0.7119^K$, namely $1.1\cdot10^{-3}$ at $K=20$,
$1.5\cdot10^{-4}$ at $K=26$ and $1.9\cdot10^{-5}$ at $K=32$ --- against the reported concentration, where
the top $10^{-3}$ of Lebesgue measure carries $55\%$, $80\%$ and $95\%$ at those same three values, and
the top $10^{-5}$ carries $31\%$ at $K=32$. The predicted carrier crosses $10^{-3}$ exactly where the
reported share stops being partial, and at $K=32$ sits at about twice $10^{-5}$, which is what a $31\%$
share of it wants. Four independently produced quantities, one of them derived from the other three,
agreeing at three depths.

It is worth seeing what a peak is, because it says where the counting programme leads. A factor is near
its maximum when $t(3/2)^m/3$ is near an integer, so a tall peak is a $t$ at which many of those
quantities are simultaneously small. At the reported top peak $t=1.869287$, fourteen of the first
twenty-eight are within $0.15$ of an integer, against eight at the generic point $t=\tfrac32$; and
$1.4046^K=2^{0.4903K}$ is exactly what half the factors contributing a two apiece would give. The
threshold is arbitrary and the reading is indicative rather than established --- a second listed peak has
only five of twenty-eight within $0.15$ and is still tall --- but the shape is clear enough. Counting the
peaks means counting the $t\in[1,2]$ for which $\|t(3/2)^m/3\|$ is small for a positive proportion of
$m\le K$, and that is a question about the orbit of $(3/2)^m$ modulo one: Mahler's problem, where even
density is unknown. The programme is concrete, and it is concrete about something hard.

One consequence deserves emphasis because it retires a control used throughout this note. Writing the
factors as $1+c\cos$, the level and volatility of $A_K$ move smoothly through $c=1$. At
$c=0.6,0.8,0.9,0.95,0.98,1$ the means are $0.498$, $0.250$, $0.140$, $0.092$, $0.067$, $0.052$ and the
increment standard deviations $0.0035$, $0.017$, $0.036$, $0.055$, $0.073$, $0.090$. The fair coin is not a phase
transition for this quantity, so a proof uniform in $c$ would be correct, and the biased-coin
substitution --- decisive elsewhere here --- carries no information about $A_K$ at all.

The limitation is worth stating in full, because it is permanent. Boundedness and logarithmic growth
cannot be told apart here by computation. Over $K=6$ to $40$ the quantity $\log K$ spans a factor
$2.06$, and the observed oscillation spans $2.12$: the data has exactly the amplitude a logarithm would
have. Separating the two would need $K$ of order $1600$, so $R$ of order $10^{280}$. Every numerical
statement in this note about $A_K$ staying bounded is therefore consistent with $A_K\sim c\log K$, which
would sink the route, and no amount of further computing will change that. The freedom to move the window is real but not
needed: splitting the octave into $m$ blocks and letting each window slide within its block, the worst
block's best constant is $0.09$ at $m=4$ and $0.21$ at $m=16$, both flat in $K$ across $K=10$ to $30$,
so a cover with a bounded constant exists --- but the direct position is already better than any of them.

\paragraph{No window can remove that growth.} Since only an upper bound on $E(R)$ is needed and
$F_K\ge0$, the sharp octave indicator may be replaced by any nonnegative majorant, and a smooth one has
a kernel decaying like $|u|^{-p}$ instead of $|u|^{-1}$. It is natural to hope the geometric growth is
an artifact of the sharp window. It is not, and the reason is visible without choosing a window. Let
$\mathrm{Cum}(u)=\sum_{0<m\le u2^K}c(m)$ be the weight the kernel meets inside $|u|$. Exhaustive
enumeration gives $\mathrm{Cum}(u)(4/3)^{-K}$ equal to
$0.15156,0.17050,0.18577,0.19095,0.19514$ at $u=\tfrac14$ and
$0.74250,0.77842,0.78679,0.79493,0.79931$ at $u=1$, for $K=6,8,10,12,14$: a fixed profile, with the
successive ratios converging to $(4/3)^2=1.7778$ at every $u$. The profile is moreover linear in $u$,
so the weight density is flat across the block at rate $(4/3)^K$ per unit of $u$, and the weight already
grows like $(4/3)^K$ inside $|m|\le\tfrac142^K$ alone. Consequently $\sum_mc(m)|G(m/2^K)|\asymp
(4/3)^K\int_0^2|G|$ for any kernel at all: fast decay does not help, and neither does vanishing to high
order at the origin, since the mass is not concentrated there. The only escape would be a kernel whose
$L^1$ norm on the block vanishes, and such a function does not majorise the indicator.

\begin{sloppypar}
What that closes, and what it does not, needs stating carefully, since the two are easily conflated.
The absolute sum bounds $\sup_{k_0}|\mathrm{Err}|$, and that supremum genuinely does grow
like $(4/3)^K$ --- so no term-by-term absolute bound on a nonnegative majorant can work, at any
$\lam\in(1/2,1)$. But the lane needs one window, $k_0=R$, and there the quantity is flat: computing
$\mathrm{Err}/L$ directly from $\sum_{k\in W}F_K(k)$, with no character expansion at all, it stays
between $-0.916$ and $-0.959$ at every even $K$ from $6$ to $30$ --- $-0.936$ at $K=6$, $-0.916$ at
$K=12$, $-0.955$ at $K=20$, $-0.957$ at $K=30$, the last on a window of $191\,751$ terms. The $(4/3)^K$ enters at the triangle
inequality and is not in the lane. So what is closed is every argument that discards the phases; an
argument that keeps them and uses the octave position is not touched by any of this.
\end{sloppypar}

That closure reads a linear profile as the signature of an $L^2$ density, so it is worth checking that
the reading discriminates rather than always returning the same answer. Run the same estimator on a
measure known to be singular: the middle-thirds Cantor measure, of correlation dimension
$\log2/\log3=0.6309$. At radii $0.05,0.02,0.01,0.005,0.002$ the ratio $\mathbb P(|D|\le r)/r$ is
$2.398,2.409,2.411,2.413,2.414$ for $\mu$, flat to $0.7\%$ and converging to $2\|f\|_2^2=2.4126$, and
$2.86,3.87,5.54,6.64,9.30$ for the Cantor measure, growing without settling, with local exponent
$0.649,0.654,0.628,0.643,0.641$ against the predicted $0.6309$. The estimator separates them cleanly,
so the linear profile is a real signature and not an artifact. The ratio is also much the sharper of
the two tests at these radii: for $\mu$ the naive exponent $\log\mathbb P/\log r$ reads
$0.708,0.775,0.809,0.834,0.858$, still well short of $1$, because $\mathbb P\sim cr$ with $c\ne1$ makes
that estimator converge only logarithmically. A reader who tested only the exponent would conclude the
correlation integral is sublinear, and be wrong.

A sharper control is available inside the same family, and it is the one to use. Let
$Y_p=\tfrac13\sum_kB_k(2/3)^k$ with $\mathbb P(B=1)=p$. Then
$\hat\mu_p(3\xi)/\hat\mu_p(2\xi)=(1-p)+pe^{-2\pi i\xi}$, which equals $1$ at every integer, so
\eqref{eq:fe} holds for \emph{every} $p$, while for $p<0.14028$ the measure is singular, its dimension
$H(p)/\log(3/2)$ falling below $1$. The functional equation therefore carries no information about
absolute continuity by itself, and any argument that uses only \eqref{eq:fe} is refuted by substituting
$p=0.1$. That substitution is the right test for everything in this note, because it changes nothing
about the ratio, the lattice, or the period --- only the bias. Applying it to the correlation-integral
reading, at depth $K=12$ the ratio $\mathbb P(|Y_p-Y_p'|\le r)/r$ across
$r=0.05,\dots,0.002$ is
\[
  \begin{array}{lll}
    p=0.5    &2.398,2.410,2.414,2.413,2.404 &\text{flat}\\
    p=0.2113 &3.81,4.39,4.83,5.19,5.64      &\text{growing}\\
    p=0.10   &7.82,12.77,18.19,25.36,41.15  &\text{growing sharply}
  \end{array}
\]
so the reading separates $p=1/2$ from the singular case at the same ratio and the same lattice. It also
carries a warning for the second moment below: $p=0.2113$ is exactly the auxiliary bias that appears
there, and its correlation integral is visibly not linear at this depth.

The profile also names itself, which says what the growth really is. Since $c(m)=2^{K+1}\mathbb
P(V=m)$ below, $\mathrm{Cum}(u)=2^K\,\mathbb P\bigl(0<|3(Y-Y')|\le u(2/3)^K\bigr)$, so a linear
correlation integral forces $\mathrm{Cum}(u)\sim\tfrac23\|f\|_2^2\,u\,(4/3)^K$. Measured at $u=1$ the
sequence extrapolates geometrically to $0.804446$ against $\tfrac23\|f\|_2^2=0.804207$. The identity
$c(m)=2^{K+1}\mathbb P(V=m)$ itself holds exactly, to $0$ in floating point over all $3^{K+1}$ patterns
at $K=3,5,7$. So the $(4/3)^K$ is not an obstruction sitting between us and the theorem: it is the
theorem, seen in the size of the absolute sum, and a proof must cancel exactly that much and no more.

The circle closes exactly here, and it is worth seeing how. The weight is a probability: $c(m)=2^{K+1}
\mathbb P(V=m)$ where $V=\sum_i\eta_i2^i3^{K-i}$ with $\eta_i$ independent, taking $0$ with probability
$\tfrac12$ and $\pm1$ with probability $\tfrac14$. But $\eta=B-B'$ for independent fair bits, so
\[
  V/3^K=\sum_i\eta_i(2/3)^i=3(Y-Y'),\qquad Y,Y'\ \text{independent copies of}\ \mu .
\]
Checked against the exact law at $K=9$ and a Monte Carlo for $3|Y-Y'|$: the two agree to three digits at
every radius from $0.5$ down to $0.02$. So the local limit theorem the reformulation asks for is a
statement about $\mathbb P(|Y-Y'|\le r)$, the correlation integral of $\mu$, and the bound it needs is
$O(r)$ exactly, which is $\mu\in L^2$. Measuring that integral from the exact law at $K=11$, the ratio
$\mathbb P(|3(Y-Y')|\le r)/r$ is $0.7702,\,0.7969,\,0.8021,\,0.8038,\,0.8042$ at
$r=0.5,\,0.2,\,0.1,\,0.05,\,0.02$, flat until the truncation scale $(2/3)^{11}=0.0116$ is reached. For an
$L^2$ density the slope should be $\tfrac23\|f\|_2^2=\tfrac23\cdot1.206310=0.80421$, which is the
measured value to four digits: the identification, the $L^2$ norm and the correlation integral all agree.
A finer scan, $400$ radii spaced logarithmically at $K=12$, gives $\mathbb P(|Y-Y'|\le r)/r$ with maximum
$2.4154$ against the predicted $2\|f\|_2^2=2.4126$, and over $r<0.1$ the whole range of the ratio is
$2.3702$ to $2.4154$, an oscillation under two per cent. Repeating at $K=14$, which reaches radii a factor of two finer, gives the same
window $2.3702$ to $2.4139$. There is no sparse sequence of radii on which the ratio grows; whatever a
counterexample would look like, it does not show here. The reformulation is therefore genuinely a change of coordinates
and nothing more; but it is a change of coordinates in which the cancellation has a visible cause rather
than being a coincidence among $2^{K+1}$ phases.

The correlation integral has its own exact recursion, which is worth recording because it says where the
linear behaviour comes from. Since $3(Y-Y')=\eta+2(Y-Y')'$ in law, with $\eta=B-B'$ independent of the
copy, writing $C(t)=\mathbb P(|Y-Y'|\le t)$ and $M(t)=\mathbb P(|\,|Y-Y'|-\tfrac12|\le t)$ gives
\[
  C(u)=2\,C(2u/3)-\tfrac12M(u),
\]
which holds against the exact law at $K=11$ to one part in $10^4$ at every radius from $0.4$ down to
$0.03$.
Read downward it is a contraction. Writing $\kappa(u)=C(u)/u$ and $m(u)=M(u)/2u$, the same relation is
\[
  \kappa(2u/3)=\tfrac34\bigl(\kappa(u)+m(u)\bigr),
\]
verified to $10^{-4}$ at every radius from $0.4$ to $0.02$. So passing to smaller scales contracts by
$3/4$ with $m$ as source, and $\kappa$ stays bounded if and only if $m$ does; the fixed point is
$\kappa=3m$, and measured $\kappa\to2.4126$ against $m\to0.8042$, which is that relation to five digits.
That fixed point is worth a second look, because it is the same statement as \eqref{eq:half} reached from
a different direction, and it is what pins the factor $\tfrac12$ in the recursion. If $A$ is the
autocorrelation density then $C(u)\sim2A(0)u$ and $M(u)\sim4A(\tfrac12)u$, so the linear limit of the
recursion reads $2A(0)=\tfrac43\cdot2A(0)-\tfrac12\cdot4A(\tfrac12)$, that is $A(0)=3A(\tfrac12)$, which is
\eqref{eq:half} exactly. Clearing denominators the law at depth $K$ lives on the lattice $3^{-(K+1)}$ and
can be built by integer shift-adds; at $K=13$ this gives $C(u)/M(u)=1.500107$ at $u=0.005$, with
$C(u)/u\to2.4132$ and $M(u)/u\to1.6087$ against the predicted $2\|f\|_2^2=2.412620$ and
$\tfrac43\|f\|_2^2=1.608413$. The recursion written with $M$ in place of $\tfrac12M$ would force
$C/M\to3$ and $A(0)=6A(\tfrac12)$, so the identity excludes it by a factor of two.
The linear behaviour of $C$ is therefore driven entirely by the mass near the half-integer translate, and
the question has moved rather than closed. The source itself satisfies a recursion involving further translates, and one can
say exactly how badly the resulting system fails to close. Writing $C_a(u)=\mathbb P(|Y-Y'-a|\le u)$, the
same computation gives
\[
  C_a(u)=\sum_{\eta}p_\eta\,C_{(3a-\eta)/2}(3u/2),
\]
so the translate moves by $a\mapsto(3a-\eta)/2$, which \emph{expands} by $3/2$. Starting from $a=0$ the
reachable translates inside $[-1,1]$ number $3,7,15,31,\dots,2^{n+1}-1$ after $n$ steps, and they are
exactly the dyadic rationals $j/2^n$: the orbit generates a dense set. The system over translates is
therefore infinite-dimensional and dense, not finite, which is the structural reason the
transfer-operator picture over translates does not close. Controlling $C$ at the origin requires
controlling $\mu$ near every dyadic point at once.

The system has a name once one passes to densities. Writing $A(a)=\int f(x)f(x+a)\,dx$ for the
autocorrelation, $C_a(u)\sim2uA(a)$ as $u\to0$, and the translate recursion becomes a functional equation
for $A$ alone:
\begin{equation}\label{eq:auto}
  A(a)=\tfrac32\sum_\eta p_\eta\,A\!\Bigl(\frac{3a-\eta}{2}\Bigr).
\end{equation}
Measured at $u=0.02$ and $K=13$ it holds to $7\cdot10^{-4}$ at every $a$ tried, and $A(0)=1.2043$ against
$\|f\|_2^2=1.206310$. At $a=0$ the two outer terms coincide by symmetry and \eqref{eq:auto} collapses to
$A(0)=3A(\tfrac12)$, which is exactly the half-shift identity \eqref{eq:half}: measured
$A(0)/A(\tfrac12)=2.9930$. So \eqref{eq:half}, the family of relations after it, and the whole translate
system are one equation read at different points, and the question is whether \eqref{eq:auto} admits a
bounded solution, which is $f\in L^2$.

Discretising the operator behind \eqref{eq:auto} on grids of $301$ and $601$ points gives a top eigenvalue
of $0.99996$ with the next moduli near $0.76$ and $0.80$: eigenvalue one with a gap of about $0.2$, which
is why the numerics converge so cleanly. It settles nothing, and the reason is worth stating, because it
is the same trap as assuming a density: a finite grid imposes boundedness by construction, and the
discretised operator is a nonnegative matrix whose Perron eigenvector exists whatever the continuum does.
What the computation does say is that if a solution exists in a space where the discretisation converges,
it is unique and stable.

The eigenvalue is exactly one, not approximately. Applying the operator to the constant and integrating,
the three branches contribute $\tfrac12,\tfrac14,\tfrac14$ after accounting for the mass that leaves
$[-1,1]$, and the total is exactly $1$; the $0.99996$ above is grid error. So the constant functional is a
left eigenvector at eigenvalue one, and $1$ lies in the spectrum unconditionally. This sharpens the
statement of the problem as far as it will go: the autocorrelation \emph{measure} of $\mu$ always
satisfies \eqref{eq:auto}, whatever $\mu$ is. What is open is only whether that measure has a bounded
density.

Iterating \eqref{eq:auto} on a grid of $4001$ points from the constant function converges in $200$ steps
to a residual of $2\cdot10^{-9}$, and the normalised limit is $A(0)=1.2063$, $A(\tfrac14)=0.9225$,
$A(\tfrac12)=0.4021$, $A(\tfrac34)=0.0749$, against the values $1.204,0.922,0.402,0.075$ measured from the
law directly. In particular it returns $\|f\|_2^2=1.2063$, agreeing with the meet-in-the-middle value
$1.206310$ to five digits: a fifth independent route to that number. The same caveat applies as above,
that the grid supplies the boundedness the continuum question is about.

That caveat can be tested, and it is the sharpest numerical evidence in this note. If the continuum fixed
point were a singular measure, refining the grid would make its peak grow: the value at the origin would
behave like $N^{1-d}$ for a local dimension $d<1$. Running the iteration at $N=501,1001,2001,4001,8001,
16001$, a refinement by a factor of $32$, gives $A(0)/\!\int\!A$ equal to
$1.206361,\,1.206323,\,1.206316,\,1.206312,\,1.206311,\,1.206311$, and the supremum equals the value at the
origin throughout. It converges, to the meet-in-the-middle value $1.206310$ to six digits, rather than
growing. This is not a proof, since the interpolation smooths at scale $1/N$ and a sufficiently slowly
growing singularity would survive such a test, but it is the test one would run to look for one, and
nothing appears.

Counting the branching explains the constant that appears everywhere in this note. Of the three children
of each translate, two stay inside $[-1,1]$ and one escapes, so the active set doubles at every step while
the scale contracts by $2/3$; the product is $2\cdot\tfrac23=\tfrac43$. That is the same $4/3$ by which the
size-only bound overshoots, the same $4/3$ separating $\tfrac12$ from $\tfrac23$ in the Riesz exponent,
and the same $4/3$ in the allowance a window bound would be given. The routes are not failing for
unrelated reasons: they are all reading one branching-times-contraction ratio that sits on the wrong side
of one.

\paragraph{A test function returns the autocorrelation.} The cleanest statement of why the majorant
idea cannot work needs no truncation, no window shape and no counting. Since the $|\hat\mu(k)|^2$ are
exactly the Fourier coefficients of the autocorrelation density $A$, for any $M:\mathbb Z\to\mathbb R$
with enough decay,
\begin{equation}\label{eq:testfn}
  \sum_kM(k)|\hat\mu(k)|^2=\int_0^1A(t)\,\check M(t)\,dt,\qquad \check M(t)=\sum_kM(k)e(kt).
\end{equation}
So every nonnegative test function dominating $1_W$ returns an integral of $A$ against its own
transform: it does not bound the correlation integral, it hands it back. Checked at $R=4096$ with an
explicit raised-cosine majorant, the two sides agree to a relative $2\cdot10^{-12}$, and the series for
$A$ returns $A(0)=1.206310$, a third independent route to $\|f\|_2^2$.

The earlier version of this paragraph set the collapse threshold at $Q\asymp(3/2)^K$ and observed that
$\min_{\eps\ne0}|\nu_\eps|=1$ for every $K\le22$, flat, while $(3/2)^K$ runs to $7482$; that verdict was
right but its threshold was too generous, since efficiency forces $Q=dP\ge P/L=3\cdot2^K$, larger by
$3(4/3)^K$. Dissociateness says no combination vanishes; it says nothing about how near zero one can
come. One caution on the other side, since it was proposed and does not hold: there is no per-octave
floor of $2\|f\|_2^2$ for majorant bounds. That would follow from $\int_{-d}^d\check M=M(0)\ge1$, but an
octave window does not contain the origin, so domination gives no lower bound at $k=0$ --- measured,
$M(0)=0$ for the majorant above, which nonetheless costs only a factor $1.59$ over $E(R)$ itself.
Majorants are not lossy here; they are simply circular.

\paragraph{A size-only bound cannot work.} Bounding $\mathrm{Err}$ by absolute values exceeds the
requirement by a factor growing like $(4/3)^K$, matching $\sum_\eps c_\eps=2^{K+1}$ against a requirement
of order $(3/2)^K$.

\paragraph{How much cancellation there is.} Computing the signed $\mathrm{Err}$ exactly, two ways that
share no code and agree to six digits:

\begin{center}
\begin{tabular}{lccccccc}
\toprule
$K$ & 3 & 5 & 7 & 9 & 11 & 12 & 13\\
\midrule
$\mathrm{Err}/L$ & $-0.867$ & $-0.978$ & $-0.949$ & $-0.936$ & $-0.925$ & $-0.917$ & $-0.938$\\
bound$/L$ & $10.8$ & $26.0$ & $63.4$ & $140.2$ & $295.5$ & $422.7$ & $606$\\
\bottomrule
\end{tabular}
\end{center}

The quantity a proof must bound above is flat at about $-0.93$, and negative, so the window sum falls
below its main term. Over the eleven values its mean is $-0.926$ with standard deviation $0.038$, and a linear
fit in $K$ has slope $-0.0046$ against a standard error of $0.0036$, so the trend is indistinguishable from flat;
the window-to-period ratio averages $0.074$. The only available bound on it is off by a factor $646$ by $K=13$ and grows
geometrically. Both methods are limited by the period and stop at $K=13$.

\paragraph{The octave form is an equivalence, not a reduction.} (The symbol here is local: $\widetilde I_K$ denotes a different normalisation from the octave integral $I_K$ of the main text.) With $\widetilde I_K=\int_1^{3/2}P_K^2$, the bound
$\widetilde I_K\le Cq^K$ with $q<2/3$ implies $\mu\in L^2$, and $\mu\in L^2$ implies $\widetilde I_K\le C(2/3)^K$ by a Fej\'er
minorant. The substitution $y=(3/2)^Kx$ makes this visible in one line:
\[
  \widetilde I_K=(2/3)^K\int_{(3/2)^K}^{(3/2)^{K+1}}\prod_{i=0}^{K}\cos^2(\pi(2/3)^iy)\,dy,
\]
so $\widetilde I_K$ is $(2/3)^K$ times a Fourier energy over one octave. Checked to six figures at
$K=6,10,14,18$.

\paragraph{The spectral gap lands exactly on the threshold.} For the self-similar variable
$D=\eps/3+(2/3)D'$ the composition operator $\mathcal C\varphi(t)=\sum_\eps p_\eps\varphi(\eps/3+(2/3)t)$
is compact on analytic functions on a disc of radius $>1$, since each branch maps it strictly inside
itself. Its spectrum is exactly $\{(2/3)^n\}\cup\{0\}$: all three branches share the derivative $2/3$, so
$\mathcal C(t^n)=(2/3)^nt^n+\text{lower order}$ and the operator is triangular in the monomial basis. A
Chebyshev discretisation at $24,40,60,80$ nodes returns $1,0.666667,0.444444,0.296296,\dots$ The second
eigenvalue is $2/3$ exactly, so no argument whose only input is this gap can produce the strict
inequality.

\paragraph{Phase-blind arguments cannot work.} The relevant small-ball quantity has infimum exactly
$2/3$ by pigeonhole and Fekete, and attaining that infimum is \emph{equivalent} to $\mu\in L^2$. Every
Frostman-type or small-ball argument therefore yields a dimension bound and provably never $L^2$.

\section{Numerical record}

All figures come from methods sharing no code: exact prefix densities at depths $20,24,26$; a
stationary-equation solver on grids of $1.8$, $3.5$ and $7.1$ million cells; a twenty-million-draw Monte
Carlo; the closed cosine product; and exact integer counts.

\begin{center}
\small
\begin{tabular}{ll}
\toprule
fibre deficits & $j_0=\log(9/8)=0.117783$, then $0.002024$, $0.000584$,\\
& $4.95\cdot10^{-5}$, $1.35\cdot10^{-5}$, $1.30\cdot10^{-5}$\\
their sum & $0.120467$ nats, equal to $\int f\log f$ to fourteen digits\\
$L^2$ norm & $\int f^2=1.206310$ and $\sum_{k\ge1}|\hat\mu(k)|^2=0.103155$,\\
& the two sides meeting as $1+2S=\|f\|_2^2$\\
octave energy & $E(R)\approx R^{-0.78}$ over octaves $4$ to $18$, stable across sub-windows\\
octave maxima & $\max|\hat\mu|\,R^{1/2}$ within $0.034$--$0.073$ over octaves $14$ to $24$,\\
& no monotone trend, though the last six drift mildly upward;\\
& the truncated $P_K$ grows and is a different object\\
weight profile & $\mathrm{Cum}(u)(4/3)^{-K}\to\tfrac23\|f\|_2^2\,u$; measured $0.804446$ at $u=1$\\
endpoint exponent & $f(x)\approx x^{0.71}$, matching $\log2/\log(3/2)-1=0.7095$\\
window constant & $\int_0^1\exp(\sum\varphi)$ over $(4/3)^K$ equal to $0.402028$ at\\
& $K=16,\dots,24$, six digits flat; the remainder bounded\\
growth constant & $6\int_0^MQ^2=1.205739,1.206198,1.206295,1.206307,1.206310$\\
& at $M=10,\dots,10^5$, against $1+2S=1.206310$\\
half-shift, second route & $\int\cos(3\pi v)Q^2/\int Q^2=0.33333349$ against $1/3$, i.e.\\
& $\rho(0)=3\rho(3)$, which is $A(0)=3A(1/2)$ in $\nu$'s coordinates\\
tail cross-ratio & $0.032,-0.039,-0.008,-0.041$ from $M=10$ to $10^4$; the $1/3$ sits\\
& entirely in $u<10$, which carries $99.95\%$ of the mass\\
renewal correlation & mean $0.978$, s.d.\ $0.087$, s.e.\ $0.020$, slope $+0.001$ per step\\
& over $n=8,\dots,27$; exponent $0.775$ against $E(R)\approx R^{-0.78}$\\
comparability & $C_K=(3/4)^KI_K/(3B_K)$ runs $119.0$ to $244.3$ over $K=8,\dots,35$,\\
& median $187.1$, factor $2.05$; bounded, not constant\\
octave rate & per-step rate of $I_K$ over $K=8,\dots,35$: arithmetic $1.0032$,\\
& geometric $0.9973$, OLS $0.9829$; sign undetermined, Jensen gap $0.0061$\\
escape rate, $1/\varphi$ & exact $0.6202985760$ four ways; float filter converges to\\
& $0.618014\pm0.000039=\lam$; smoothing gives $\lam$ at every resolution\\
lines near the target & $\#\{|3^{n+1}-S_\eps|\le D2^{\,n+1}\}=D\,2^{\,n}$ exactly, $n=4,\dots,16$,\\
& matching uniform spreading; nearest line at distance $2^{-(n+1)}$\\
one-step ratios & $0.019,\ 0.289,\ 0.085,\ 0.273,\ 0.966,\ 0.007,\ 0.119,\ 0.736,\ 0.533,\ 0.187$\\
\bottomrule
\end{tabular}
\end{center}

The fifth of those ratios deserves a line: the contraction of the fibre bias is not uniform in $n$, and
any argument assuming it is will fail at level four.

\end{document}
